% concatenated from arxiv.tex, packages.tex
\documentclass{article}
\PassOptionsToPackage{numbers, compress}{natbib}
\usepackage[preprint]{neurips_2026}
\usepackage[utf8]{inputenc}
\usepackage[T1]{fontenc}
\usepackage{hyperref}
\usepackage{url}
\usepackage{booktabs}
\usepackage{amsfonts}
\usepackage{nicefrac}
\usepackage{microtype}
\usepackage{xcolor}
\usepackage{subcaption}
\usepackage{titletoc}
\bibliographystyle{plainnat}
\usepackage{amsmath,mathtools,bm}
\usepackage{amsthm}
\usepackage[T1]{fontenc}
\usepackage{microtype}
\usepackage{enumitem}
\usepackage{booktabs,tabularx,array, makecell}
\hypersetup{colorlinks=true, linkcolor=blue!50!black, urlcolor=blue!50!black, citecolor=blue!50!black}
\newtheorem{theorem}{Theorem}
\newtheorem{proposition}[theorem]{Proposition}
\newtheorem{lemma}[theorem]{Lemma}
\newtheorem{corollary}[theorem]{Corollary}
\newtheorem{remark}[theorem]{Remark}
\newtheorem{definition}[theorem]{Definition}

\newcommand{\X}{\mathcal{S}}
\newcommand{\A}{\mathcal{A}}
\newcommand{\Fcal}{\mathcal{F}}
\newcommand{\Pcal}{\mathcal{P}}
\newcommand{\R}{\mathbb{R}}
\newcommand{\E}{\mathbb{E}}
\newcommand{\Prob}{\mathbb{P}}
\newcommand{\Tc}{T^{\pi}}
\newcommand{\norm}[1]{\left\lVert #1 \right\rVert}
\newcommand{\ip}[2]{\left\langle #1, #2 \right\rangle}
\newcommand{\Ustar}{U^{\star}}
\newcommand{\Vstar}{V^{\star}}


\title{Quotient-Categorical Representations for Bellman-Compatible Average-Reward Distributional Reinforcement Learning}

\author{Ege C.~Kaya, Aliasghar~Pourghani, Vijay~Gupta, Abolfazl~Hashemi  \\
  Elmore Family School of Electrical and Computer Engineering\\
  Purdue University\\
  West Lafayette, IN 47906 \\
  \texttt{kayae@purdue.edu, apourgha@purdue.edu, gupta869@purdue.edu, abolfazl@purdue.edu}}

\begin{document}
\maketitle

\begin{abstract}
Average-reward reinforcement learning requires estimating the gain and the bias, which is defined only up to an additive constant. This makes direct distributional analogues ill-posed on the real line. We introduce a quotient-space formulation in which state-indexed bias laws are identified up to a common translation, together with a categorical parameterization that respects this symmetry. On this quotient-categorical space, we define a projected average-reward distributional operator and show that it is well-defined, non-expansive in a coordinate Cram\'er metric, and admits fixed points. We then study sampled recursions whose mean-field maps are asynchronous relaxations of this operator. In an idealized centered-reward setting, a one-state temporal-difference update enjoys almost sure convergence together with finite-iteration residual bounds under both i.i.d.\ and Markovian sampling. When the gain is unknown, we augment the recursion with an online gain estimator, and prove non-expansiveness and Markovian convergence of the resulting coupled scheme. Finally, we show that synchronous exact updates are gain-independent at the quotient-law level, isolating a structural contrast between ideal quotient distributions and practical fixed-grid categorical representations.
\end{abstract}

\section{Introduction}

Distributional reinforcement learning (DRL) is well developed in the discounted setting, where the standard contraction arguments and Bellman-compatible return distributions support both analysis and algorithms \citep{bellemare2017distributional, bellemare2023distributional, rowland2018analysis, dabney2018distributional, dabney2018implicit}. The average-reward setting is structurally different. When the discounting is removed, the contraction geometry is lost, and fixed-policy evaluation is no longer described by a discounted return. Instead, average-reward evaluation revolves around two objects: the \emph{gain}, which is the long-run average reward, and the \emph{bias}, which solves the average-reward analogue of the Bellman equation, namely the Poisson equation. The bias is the object that distinguishes states and policies with the same gain, captures transient differential performance pre-stationarity, and underlies average-reward evaluation and policy improvement \citep{puterman, sutton1998reinforcement, mahadevan1996average}. However, it is defined only up to an additive constant. This creates an obstruction for distributional modeling: a distributional bias cannot be represented as an ordinary state-indexed family of laws on $\R$ with a fixed origin, because translating every state's bias by the same constant should not change the object being learned.

We therefore model state-indexed bias laws modulo common translations. The resulting quotient viewpoint preserves the additive-constant symmetry of the Poisson equation, and our categorical construction turns this quotient object into a finite-dimensional representation. This makes it possible to define projected Bellman operators for the distributional bias, compare categorical iterates in a coordinate Cram\'er metric, and write temporal-difference (TD) algorithms whose fixed points are Bellman-compatible projected quotient-categorical bias laws. Our objective is complementary to the average-reward distributional criterion studied, e.g., by \citet{rojas2026differential}. That line of work refines the average-reward criterion itself through limiting per-step reward distributions. We instead refine the bias: the Bellman-compatible object that records state-dependent differential performance. Thus the goal is not merely to describe long-run reward statistics, but to construct a distributional analogue of the scalar bias that can be used in average-reward dynamic-programming and TD-learning arguments.

The remaining challenge is algorithmic. The projected maps we obtain are non-expansive rather than contractive, so discounted DRL contraction arguments do not apply. We show that the sampled categorical updates can be written as TD / stochastic approximation (SA) recursions in the classical sense \citep{robbins1951stochastic, sutton1988learning, tsitsiklis1994asynchronous, borkar1998asynchronous, borkar2008stochastic}, identify their mean-field maps, and prove that these maps are non-expansive and share fixed points with the projected bias operator. This connects the quotient-categorical bias recursion to recent SA results for non-expansive maps \citep{bravo2024stochastic, blaser2026asymptotic}.

\noindent\textbf{Contributions.} We put average-reward distributional bias estimation on firm foundations by defining the correct quotient-categorical target and proving convergence guarantees for TD-style categorical recursions. Our main contributions are as follows.

\begin{enumerate}[leftmargin=*,labelsep=0.2em, align=left, itemsep=2pt]
\item We establish that the quotient space of state-indexed law families by common translation is the natural space for a distributional analogue of the scalar bias.
\item We define an exact quotient-space operator and a quotient-categorical projection algorithm, proving in particular that the projected operator is well-defined, non-expansive, and admits a fixed point.
\item We analyze a centered-reward TD recursion for the quotient-categorical bias distribution, whose mean-field maps are non-expansive asynchronous relaxations of the projected operator, with a.s.\ convergence and finite-iteration residual bounds.
\item We introduce a practical online gain-estimated one-state TD recursion, which is non-expansive and admits a mean-field map with correct fixed points.
\end{enumerate}

\section{Related work}
DRL is built around Bellman-compatible return distributions in discounted MDPs, with categorical and quantile parameterizations providing the dominant practical representations \citep{bellemare2017distributional, rowland2018analysis, dabney2018distributional, dabney2018implicit}. Early theory studied the geometry, statistics, and relative advantages of categorical and quantile DRL \citep{rowland2018analysis, lyle2019comparative, rowland2019statistics}. Recent work has made many advancements by sharpening the statistical case for learning return distributions \citep{wang2023benefits, wang2024more}, providing non-asymptotic efficiency of distributional TD \citep{peng2024statistical, peng2025finite}, near-minimax guarantees in the generative-model regime \citep{rowland2024near}, establishing multivariate distributional dynamic programming \citep{wiltzer2024foundations}, offline distributional policy evaluation guarantees \citep{wu2023distributional}, and convergence analyses for KL-trained categorical updates \citep{kastner2025categorical}. In the average-reward setting, however, most work remains scalar, from early R-learning and average-reward foundations to modern off-policy, robust, and function-approximation results \citep{schwartz1993reinforcement,abounadi2001learning,mahadevan1996average, wan2021learning, wang2023model, vakili2024kernel, zurek2024span}. The closest prior average-reward DRL work is \citep{rojas2026differential}, where the authors study long-run per-step reward and differential return distributions via quantile representations. As opposed to these works, we seek a Bellman-compatible distributional analogue of the bias, which makes the quotient symmetry induced by additive-constant non-uniqueness central rather than incidental. On the algorithmic side, our finite-iteration analysis draws on fixed-point methods for non-expansive maps \citep{krasnosel1955two, mann1953mean, kim2007robustness,bravo2019rates,bravo2024stochastic,blaser2026asymptotic} and complements classical TD / SA analyses for contractive or linear settings \citep{dayan1994td,tsitsiklis1996analysis,bhandari2018finite,dalal2018finite,srikant2019finite,patil2023finite,qu2020finite,de2020fixed,chen2020finite,chen2022finite,chen2024lyapunov}.
\section{Preliminaries}\label{sec:prelim}

We work with a finite average-reward MDP \citep{puterman} $(\X, \A, R, P)$, where $\X = \{1,\dots,m\}$ is the finite set of states and $\A$ is the finite set of actions, $R: \X \times \A \to [0, 1]$ is the deterministic reward function and $P:\X \times \A \to \Delta(\X)$ is the transition kernel. A policy $\pi$ is fixed throughout, and we focus on its evaluation. With $\pi$ fixed, the induced process is a Markov reward process with transition matrix $P^\pi$, where $P^\pi_{ij} = \sum_{a\in\A}P(j \mid i, a) \,\pi(a \mid i)$ \citep{sutton1998reinforcement}. Throughout the paper, we assume that the induced Markov chain $\{S_t\}$ is irreducible and aperiodic, and write $\mu$ for its unique stationary distribution. We denote by $R_{ij}$ the finite-valued random one-step reward conditioned on $(S_t = i, S_{t+1} = j)$. We define the expected one-step reward vector as
\begin{equation}\label{eq:average-reward}
r^\pi_i := \sum_{j\in \X} P^\pi_{ij} \E[R_{ij}],\quad \text{for all } i \in \X.
\end{equation}

The (scalar) \emph{gain} or \emph{average reward} is then
\begin{equation}
\bar r^\pi := \sum_{i\in \X} \mu_i \sum_{j\in\X}P^\pi_{ij} \E[R_{ij}] = \sum_{i \in \X} \mu_i r^\pi_i.
\end{equation}

The \emph{bias} or \emph{relative value function} \citep{puterman} $v^\pi \in \R^m$ is any solution of the Poisson equation \citep{Meyn_Tweedie_Glynn_2009}
\begin{equation}\label{eq:poisson}
v^\pi = r^\pi - \bar r^\pi \mathbf 1 + P^\pi v^\pi.
\end{equation}
The Poisson equation is, in some sense, the analogue of the Bellman equation \citep{bellman1966dynamic} in the average-reward setting, where the discount factor has been set to $1$ and the rewards have been centered. 

The Poisson equation can equivalently be written as $(I- P^\pi) v^\pi = r^\pi - \bar r^\pi \mathbf 1$. Since the null space of $(I-P^\pi)$ contains any constant vector $c\mathbf 1 \in \R^m$, the solutions of \eqref{eq:poisson} are unique up to addition of a common constant. This scalar translation-invariance property of solutions of \eqref{eq:poisson} provides the intuition for the quotient-categorical formulation constructed in the rest of the work.

\section{Quotient-categorical bias estimation}\label{sec:quotient}

In this work, we seek to estimate the distribution of the bias. 
As motivated by the discussion in Section~\ref{sec:prelim}, our first objective is to recover an analogue of the translation-invariance property in the distributional setting. Let $\mathcal F^\X := \{\eta = (\eta_i)_{i \in \X}: \eta_i \in \mathcal P_2(\R)\}$ denote the set of state-indexed collections of probability laws on $\R$ with finite second moment. For $c \in \R$, let $\tau_c(x) := x+c$, and define \emph{common-translation} on $\mathcal F^\X$ by $((\tau_c)_\# \eta)_i := (\tau_c)_\# \eta_i$, where $\#$ denotes pushforward. We also define the equivalence relation $\approx$ on any pair $\eta, \zeta \in \mathcal F^\X$ by
\begin{equation}
\eta \approx \zeta \iff \exists c\in\R \text{ such that } \zeta = (\tau_c)_\#\eta,
\end{equation}
and denote the quotient space by $\mathcal F^\X / \approx$. Let us write $B_{ij} := R_{ij} - \bar r^\pi$ for the centered one-step reward, and $\nu_{ij} := \mathrm{Law}(B_{ij})$. We now define the average-reward distributional operator.
\begin{definition}
The \emph{average-reward distributional operator} $\mathcal T: \mathcal F^\X \to \mathcal F^\X$ is defined componentwise:
\begin{equation}\label{eq:avg-operator}
(\mathcal T \eta)_i := \sum _{j\in \X} P_{ij}\, (\nu_{ij}\,\ast\,\eta_j), \quad \text{for all } i \in \X.
\end{equation}
\end{definition}
\begin{proposition}\label{prop:T}
For every $\eta \in \mathcal F^\X$ and $c \in \R$,
\begin{equation}\label{eq:commute}
\mathcal T\bigl((\tau_c)_\# \eta \bigr) = (\tau_c)_\# \mathcal T\eta.
\end{equation}
Consequently, $\mathcal T$ induces a well-defined operator on the quotient space $\mathcal F^\X / \approx$:
\begin{equation}
\overline {\mathcal T}: \mathcal F^\X / \approx\; \to \mathcal F^\X / \approx, \quad \overline{\mathcal T}\bigl([\eta]\bigr) :=[\mathcal T \eta].
\end{equation}
\end{proposition}

The quotient formulation is instructive in seeing how the gain is not part of the quotient target itself. If we center rewards by an arbitrary constant $g \in \R$ instead of $\bar r^\pi$, the resulting state-indexed collection of laws differs only by a common translation. The following proposition formalizes this.
\begin{proposition}\label{prop:tg}
Let $g \in \R$. Define the \emph{$g$-centered average-reward distributional operator}
\begin{equation}
(\mathcal T_g \eta)_i := \sum_{j\in\X} P_{ij}(\nu_{ij}^{(g)}\,\ast\, \eta_j),\quad\text{where } \nu_{ij}^{(g)}:= \mathrm{Law}(R_{ij} -g).
\end{equation}
Then, for any $g, g' \in \R$ and $\eta \in \mathcal F^\X$, $[\mathcal T_g \eta] = [\mathcal T_{g'} \eta]$ in $\mathcal F^\X / \approx$.
\end{proposition}

Guided by these observations, we introduce a categorical representation class which respects the translation-invariance property of the bias law. Let $\Theta = \{ \theta_1 < \cdots < \theta_d\} \subset \R$ be an ordered support with constant stride $\Delta:=\theta_{k+1} - \theta_k >0$ for $k=1,\dots, d-1$.\footnote{It is straightforward to use state-dependent supports and non-uniform stride, we forgo these for simplicity of exposition.} Let $\Delta_d$ denote the $d$-dimensional simplex. For $p \in \Delta^\X_d$ and $c \in \R$, define the representative family of laws
\begin{equation}
\eta_i^{p, c} := \sum_{k=1}^d p_{i, k}\, \delta_{\theta_k +c},\quad \text{for all } i \in \X.
\end{equation}
The state-indexed collection of coefficients $p$ therefore uniquely characterizes a quotient class of categorical distributions $[\eta^p] := [\eta^{p, c}] \in \mathcal F^\X / \approx$, independently of shift $c$. We call this a \emph{quotient-categorical} representation: categorical because
each state law is represented on the finite grid $\Theta$, and quotient because
the common shift $c$ is not part of the state, but only a choice of representative.

For each $c \in \R$, define the $c$-shifted support $\Theta_c := \{\theta +c : \theta \in \Theta\}$, and let $\Pi^{\Theta_c}_{\mathrm C}$ denote the standard linear-interpolation categorical projection \citep{bellemare2017distributional} onto $\Theta_c$, applied over all states. We will show that this projection is shift-equivariant: first shifting a distribution by $c$ and then projecting onto the $c$-shifted categorical support is equivalent to projecting it onto the unshifted support and then shifting by $c$.
\begin{lemma}[Projection shift-equivariance]\label{lem:proj-shift-equiv}
For every $\eta \in \mathcal P_2(\R)$ and $c \in \R$,
\begin{equation}
\Pi^{\Theta_c}_{\mathrm C} \bigl( (\tau_c)_\# \eta \bigr) = (\tau_c)_\# (\Pi^{\Theta}_{\mathrm C} \eta).
\end{equation}
\end{lemma}
This same fact is illustrated in Figure~\ref{fig:translation_fixed_support}. Since every representative $\eta^{p, c}$ is supported on a $c$-shifted copy $\Theta_c$ of the same support $\Theta$, its Cram\'er geometry is determined entirely by the cumulative masses over the support. For $p \in \Delta_d$, let $F_p(\theta_k) := \sum_{j=1}^k p_j$ for $ k=1,\dots,d-1$. We define, for $p, q \in \Delta_d$, the coordinate Cram\'er metric by
\begin{equation}
\ell^\Theta_{\mathrm C}(p, q)^2 := \Delta \sum_{k=1}^{d-1}\bigl(F_p(\theta_k) - F_q(\theta_k)\bigr)^2.
\end{equation}
The statewise extension of this is the so-called supremum-Cram\'er metric, where for $p, q \in \Delta^\X_d$,
\begin{equation}
\ell^\Theta_{\mathrm C, \infty}(p, q) := \max_{i \in \X}\ell^\Theta_{\mathrm C}(p_i, q_i).
\end{equation}

The next proposition shows that these coordinate metrics coincide exactly with the Cram\'er metrics on the corresponding representative families.
\begin{proposition}\label{prop:cramer}
For every $p, q \in \Delta^\X_d$ and $c \in \R$,
\begin{equation}
\ell_{\mathrm C}(\eta^{p, c}_i, \eta^{q, c}_i) = \ell^\Theta_{\mathrm C}(p_i, q_i),\quad\text{for all } i \in \X.
\end{equation}
Consequently,
\begin{equation}
\ell_{\mathrm C, \infty}(\eta^{p, c}, \eta^{q, c}) = \ell^\Theta_{\mathrm C, \infty}(p, q),\quad\text{for all } c \in \R.
\end{equation}
\end{proposition}

Proposition~\ref{prop:cramer} shows that the representative families $\eta^{p, c}$ faithfully encode the relevant Cram\'er geometry in the coefficient vectors $p \in \Delta^\X_d$, regardless of the shift $c$, as long as the same $c$ is chosen for both representatives. We now use this representation to define the projected analogue of the $g$-centered operator $\mathcal T_g$ on categorical quotient coordinates. Proposition~\ref{prop:tg} previously showed that, at the quotient level, $\mathcal T_g$ does not depend on the choice of $g$: changing $g$ only changes the representative of the output, but lands in the same equivalence class. After categorical projection, however, the representative matters, since the projection is performed relative to a fixed shifted support $\Theta_c$.

\begin{definition}
\label{defn:proj-operator}
Let $g \in \R$ and $p \in \Delta^\X_d$ and choose $c \in \R$ arbitrarily. Applying $\mathcal T_g$ to the representative family $\eta^{p, c}$ and projecting each component back onto $\Theta_c$ yields another categorical family on $\Theta_c$. We define $\mathcal G_g(p) = q \in \Delta^\X_d$ to be the unique element that satisfies
\begin{equation}
\Pi^{\Theta_c}_{\mathrm C}(\mathcal T_g \eta^{p, c})_i = \eta^{q, c}_i,\quad \text{for all } i \in \X.
\end{equation}
\end{definition}
We note that this definition is independent of the choice of the representative shift $c$. Indeed, since by definition $\eta^{p, c'} = (\tau_{c'-c})_\#\eta^{p, c}$, the same argument as in Proposition~\ref{prop:T} gives
\begin{equation}
\mathcal T_g \eta^{p, c'} = (\tau_{c' -c})_\# \mathcal T_g \eta^{p, c}.
\end{equation}

\begin{figure}[t]
    \centering
    \includegraphics[width=\textwidth]{translation_shifted_vs_fixed_support.png}
    \caption{
    The left panel illustrates Lemma~\ref{lem:proj-shift-equiv}: translating both the law and support preserves the coefficient vector $p$. 
    The right panel illustrates that projecting a translated law back onto the same support generally gives different coefficients $q\neq p$.
    }
    \label{fig:translation_fixed_support}
\end{figure}













The application of Lemma~\ref{lem:proj-shift-equiv} shows that projection onto $\Theta_{c'}$ or onto $\Theta_c$ produces categorical families with the same coefficient vector $q$. Therefore, $\mathcal G_g$ is a well-defined operator on $\Delta^\X_d$ for any $g \in \R$. Lastly, we define $\mathcal G := \mathcal G_{\bar r^\pi}$ as the ``correctly-centered'' projected operator.

We now state the main theorem of this section, establishing the non-expansiveness of $\mathcal G_g$ in the coordinate Cram\'er metric. 

\begin{theorem}[Non-expansiveness of $\mathcal G_g$]\label{thm:gg-nonexp}
For every $g \in \R$, and every $p, q \in \Delta^\X_d$,
\begin{equation}
\ell_{\mathrm C, \infty}^\Theta \bigl(\mathcal G_g(p), \mathcal G_g(q)\bigr) \le \ell_{\mathrm C, \infty}^\Theta(p, q).
\end{equation}
In particular, the correctly-centered operator $\mathcal G$ is non-expansive.
\end{theorem}
\begin{proof}[Proof sketch]
Fix a common representative shift $c$ and compare the representative families $\eta^{p, c}$ and $\eta^{q, c}$. For each state $i$, the $i$th component of $\mathcal T_g$ is a $P_{ij}$-weighted average of convolved laws $\nu_{ij}^{(g)} \, \ast\, \eta_j$. Convolution with a fixed law is non-expansive in the Cram\'er metric, and the subsequent $P_{ij}$-averaging is also non-expansive. Thus $\mathcal T_g$ is non-expansive in $\ell_{\mathrm C, \infty}$. Composing with the standard non-expansive categorical projection and then using Proposition~\ref{prop:cramer} yields the claim. 
\end{proof}

\begin{corollary}[Existence of fixed points]\label{corol:fixed-points}
Let $g \in \R$. $\mathrm{fix}(\mathcal G_g)$ is nonempty.
\end{corollary}

\section{A centered stochastic approximation recursion}\label{sec:centered}
In the previous section, we defined the $g$-centered projected operator $\mathcal G_g$ and established its non-expansiveness in the coordinate Cram\'er metric. In particular, $\mathcal G$ is a non-expansive operator with at least one fixed point. We now turn from this exact projected backup to an SA thereof, based on samples from interaction with the MDP. Customarily, SA in RL comes in the form of a TD update, where we observe only one transition at a time and thus asynchronously modify only one state block.

A one-state sample is defined to be a tuple $y = (s, b, s') \in \mathcal Y := \X \times [-1, 1] \times \X$.\footnote{Note that rewards lying in $[0, 1]$ implies that centered rewards lie in $[-1, 1]$.} For $p \in \Delta^\X_d$, we define $H(p ,y) = q \in \Delta^\X_d$, the one-sample backup map, by
\begin{equation}\label{eq:one-sample-backup}
q_u = \begin{cases}
p_u, &u\ne s,\\
\tilde q, &u=s,
\end{cases}
\end{equation}
where $\tilde q \in \Delta_d$ is the unique coefficient vector satisfying $\Pi^\Theta_{\mathrm C}\bigl( (\tau_b)_\# \eta ^{p_{s'}, 0}\bigr) = \eta^{\tilde q, 0}$. The asynchronous recursion is thus the well-known stochastic Krasnosel'ski\u{\i}--Mann (SKM) iteration \citep{krasnosel1955two, mann1953mean} 
\begin{equation}\label{eq:KM}
p_{k+1} = p_k + \alpha_k \bigl(H(p_k, Y_k) - p_k\bigr),\quad k\ge0.
\end{equation}

The next proposition rewrites the $g$-centered projected operator $\mathcal G_g$ blockwise in categorical coefficients, which is the form needed to derive asynchronous SA updates. Before the statement, we introduce a slight abuse of notation. For $u \in \Delta_d$ and $c \in \R$, we write $
\eta^{u, c} := \sum_{k=1}^d u_k \delta_{\theta_k +c}$.
In particular, $(\eta^{p, c})_i = \eta^{p_i, c}$ for $p \in \Delta^\X_d$.
\begin{proposition}\label{prop:g-block-form}
Let $g \in \R$, $p \in \Delta^\X_d$, and $i \in \X$. The $i$th block of $\mathcal G_g(p)$ is the unique vector $q_i \in \Delta_d$ such that
\begin{equation}
\eta^{q_i, 0} = \Pi^\Theta_{\mathrm C}\Bigl(\sum_{j \in \X} P_{ij}(\nu_{ij}^{(g)} \,\ast\, \eta^{p_j, 0})\Bigr).
\end{equation}
\end{proposition}
The next lemma simply states that the one-sample backup map is also a non-expansion. 
\begin{lemma}\label{lem:one-sample-nonexp}
For every $y \in \mathcal Y$ and $p, q \in \Delta^\X_d$,
\begin{equation}
\ell^\Theta_{\mathrm C, \infty}\bigl( H(p, y), H(q, y)\bigr) \le \ell^\Theta_{\mathrm C, \infty}(p, q).
\end{equation}
\end{lemma}

We will now work on creating an asynchronous relaxation of $\mathcal G$ through the use of the one-sample backup map $H$. Let $\rho \in \Delta(\X)$ such that $\rho_{\min}:=\min_{i \in \X} \rho_i >0$. We define
\begin{equation}
D_\rho := \mathrm{diag}(\rho_1, \dots, \rho_m),\qquad h_\rho(p):= p + D_\rho\bigl(\mathcal G(p) - p\bigr).
\end{equation}
$h_\rho$ is the \emph{mean-field map} of the one-sample stochastic update under state-sampling law $\rho$, in the standard SA sense. Of particular interest will be when $\rho=\mu$, the stationary distribution of the Markov chain under $\pi$. We have the following favorable results on this averaged map.

\begin{proposition}\label{prop:mean-field-fix}
Suppose $Y = (S, B, S')$ is sampled such that $S \sim \rho$ and
\begin{equation}
\mathbb P(S'= j \mid S = i) = P_{ij}, \qquad B \mid (S=i, S'=j) \sim \nu_{ij}\quad \text{for all } i, j \in \X.
\end{equation}
Then, $\E\bigl[H(p, Y)\bigr] = h_\rho(p)$ for all $p \in \Delta^\X_d$, $h_\rho$ is non-expansive in $\ell^\Theta_{\mathrm C, \infty}$, and $\mathrm{fix}(h_\rho) = \mathrm{fix}(\mathcal G)$. In particular, the same conclusions hold with $\rho = \mu$ under the standing irreducibility assumption.
\end{proposition}

The next corollary will be useful in translating residual bounds for $h_\rho$ into residual bounds for $\mathcal G$.
\begin{corollary}\label{corol:residual-conversion}
For every $p \in \Delta^\X_d$,
\begin{equation}
\rho_{\min} \ell^\Theta_{\mathrm C, \infty}\bigl(p, \mathcal G(p)\bigr) \le \ell^\Theta_{\mathrm C, \infty}\bigl(p, h_\rho(p)\bigr) \le \ell^\Theta_{\mathrm C, \infty}\bigl(p, \mathcal G(p)\bigr).
\end{equation}
\end{corollary}

We state the two main SA results for the centered-reward recursion approximating $\mathcal G$. The i.i.d.\ regime is the relevant model for settings such as replay-buffer sampling or a generative simulator, where we can approximately draw independent one-step samples from a fixed sampling law. The Markovian regime captures online asynchronous TD-learning along a single trajectory of the policy-induced chain. In both settings, the quantity we control is the mean-field residual $\ell^\Theta_{\mathrm C, \infty}\bigl(p_k, h_{\rho}(p_k)\bigr)$. This is the standard quantity in non-expansive SA \citep{mann1953mean, bravo2019rates, bravo2024stochastic, blaser2026asymptotic, kim2007robustness}, where fixed points need not be unique.

\begin{theorem}[Centered recursion i.i.d.\ convergence]\label{thm:centered-iid}
Suppose the samples $Y_k = (S_k, B_k, S'_k)$ are i.i.d., that $S_k \sim \rho$ with $\rho_{\min} >0$ and that 
\begin{equation}
\mathbb P (S'_k = j \mid S_k = i) = P_{ij}, \qquad B_k \mid (S_k =i, S'_k=j) \sim \nu_{ij}\quad \text{for all } i, j \in \X.
\end{equation}
The following two conclusions hold for the centered SKM recursion \eqref{eq:KM}. First, if $\alpha_k = (k+1)^{-a}$ with $\frac{2}{3} < a \le 1$, then $p_k$ converges a.s.\ to a random fixed point $p^\star \in \mathrm{fix}(h_\rho) = \mathrm{fix}(\mathcal G)$. Second, fix any $a_1 \in \bigl(\frac{2}{3}, 1\bigr)$ and set $\varepsilon_{a_1} := (3a_1 -2)/6$. There is an explicit two-phase step size schedule and constants $C_\mathrm{iid}, \kappa_{a_1}>0$ such that, for all $k \ge 1$,
\begin{equation}
\E\Bigl[\ell^\Theta_{\mathrm C, \infty}\bigl(p_k, \mathcal G (p_k)\bigr) \Bigr] \le \frac{C_{\mathrm{iid}}}{\rho_{\min}} (k+1)^{-1/6} \min\left\{ \kappa_{a_1}\log(k+1), (k+1)^{\varepsilon_{a_1}}\right\},
\end{equation}
where $C_{\mathrm{iid}}$ depends only on $a_1$, $\lvert \X \rvert$, and $\mathrm{diam}_{\ell^\Theta_{\mathrm C,\infty}}(\Delta^\X_d) = \sqrt{\theta_d - \theta_1}$.
The full fixed-exponent rate profile is presented in Appendix~\ref{app:B}. In particular, choosing $a = \frac{2}{3} + \epsilon$ for $\epsilon >0$ gives a.s.\ convergence with expected residual decaying at rate $O(k^{-1/6+\epsilon/2})$.
\end{theorem}
\begin{theorem}[Centered recursion Markovian convergence]\label{thm:centered-markov}
Let $(S_k)_{k\ge 0}$ be an irreducible and aperiodic Markov chain with stationary distribution $\mu$, and let $Y_k = (S_k, B_k, S_{k+1})$ be the associated centered one-sample process. The following two conclusions hold for the centered SKM recursion \eqref{eq:KM}. First, if $\alpha_k = (k+1)^{-a}$ with $\frac{4}{5} < a \le 1$, then $p_k$ converges a.s.\ to a random fixed point $p^\star \in \mathrm{fix}(h_\mu)=\mathrm{fix}(\mathcal G)$. Second, fix any $a_1 \in \bigl(\frac{4}{5}, 1\bigr)$ and set $\varepsilon_{a_1}:=(5a_1-4)/10$. There is an explicit two-phase step size schedule and constants $C_{\mathrm{mk}}, \kappa_{a_1} >0$ such that, for all $k \ge 1$, 
\begin{equation}
\E\Bigl[\ell^\Theta_{\mathrm C, \infty}\bigl(p_k, \mathcal G (p_k)\bigr) \Bigr] \le \frac{C_{\mathrm{mk}}}{\mu_{\min}}(k+1)^{-1/10} \min \left\{\kappa_{a_1}\log(k+1), (k+1)^{\varepsilon_{a_1}} \right\},
\end{equation}
where $C_{\mathrm{mk}}$ depends on $a_1$, $\mathrm{diam}_{\ell^\Theta_{\mathrm C,\infty}}(\Delta^\X_d)=\sqrt{\theta_d - \theta_1}$, and a finite-state Poisson constant for $(Y_k)$, equivalently expressable through the norm of its fundamental matrix. The full fixed-exponent rate profile is presented in Appendix~\ref{app:B}. In particular, choosing $a = \frac{4}{5}+\epsilon$ for  $\epsilon >0$ gives a.s.\ convergence with expected residual decaying at rate $O(k^{-1/10+\epsilon/2})$.
\end{theorem}
\begin{proof}[Proof sketch for Theorems~\ref{thm:centered-iid} and~\ref{thm:centered-markov}]
By Lemma~\ref{lem:one-sample-nonexp}, the one-sample backup map is non-expansive. By Proposition~\ref{prop:mean-field-fix} and Corollary~\ref{corol:fixed-points}, the mean-field map is non-expansive with a non-empty set of fixed points. Compactness of $\Delta^\X_d$ yields bounded iterates and bounded martingale noise, so the i.i.d.\ recursion falls under the bounded-noise non-expansive SA results of \citet{bravo2024stochastic}. For the Markovian regime, finiteness of $\X$ and $\A$, together with the deterministic and bounded reward function $R(s, a)$ imply that the feasible support of $Y_k$ is finite. Moreover, $(Y_k)_{k\ge0}$ is itself an irreducible and aperiodic Markov chain, inherited from $(S_k)_{k \ge 0}$, so the fixed-exponent results of \citet{blaser2026asymptotic} apply. The two-phase residual bounds are proved separately in Appendix~\ref{app:B}.
\end{proof}

The finite-iteration bounds above use a two-phase step size refinement. A fixed supercritical exponent step size schedule gives a.s.\ convergence and a residual bound free of logarithmic factors, but loses a polynomial factor relative to the last-iterate-optimal exponent. The critical exponent schedule gives the best rate up to a logarithmic factor, but pays that logarithmic factor in the transient window. The two-phase schedule gives a single finite-time bound that tracks the smaller of these two bounds. The exact schedules and proofs are given in Appendix~\ref{app:B}.

\section{A coupled recursion for when the gain is unknown}\label{sec:uncentered}

The centered recursion isolates the non-expansive structure, but it is not directly implementable unless the gain $\bar r^\pi$ is known. In practice, we observe uncentered rewards $R$, not $B:= R - \bar r^\pi$. 

\begin{proposition}\label{prop:counterexample}
In general, replacing the centered reward $B = R - \bar r^\pi$ by the raw reward $R$ in the one-sample backup \eqref{eq:one-sample-backup} does not yield an SA of $\mathcal G$. 
\end{proposition}

We now replace centered rewards by uncentered ones, and augment the quotient-categorical update with an online estimator $g$ of the gain $\bar r^\pi$, on the product space $\mathcal Z := \Delta^\X_d \times [0, 1]$. The goals of this section are threefold: first, to show that the one-sample augmented update is non-expansive. Second, to identify its mean-field map and its fixed points. Third, to relate the resulting augmented residual to the two quantities of actual interest, namely, the projected-operator residual $\ell^\Theta_{\mathrm C, \infty}\bigl(p, \mathcal G(p)\bigr)$, and the gain error $\lvert g - \bar r^\pi\rvert$. For $g \in \R$, define the gain-parameterized one-sample backup map $H_g:\Delta^\X_d \times \mathcal Y \to \Delta^\X_d$ as follows: For an uncentered sample $y=(s, r, s')$, let $H_g(p, y) = q$, where
\begin{equation}\label{eq:one-sample-backup-gain}
q_u = \begin{cases}
p_u, &u\ne s,\\
\tilde q, &u=s,
\end{cases}
\end{equation}
where $\tilde q \in \Delta_d$ is the unique coefficient vector satisfying $\Pi^\Theta_{\mathrm C}\bigl( (\tau_{r-g})_\# \eta ^{p_{s'}, 0}\bigr) = \eta^{\tilde q, 0}$. The corresponding mean-field map is defined in terms of the stationary state distribution $\mu$ as $h^{(g)}_\mu(p) := p + D_\mu(\mathcal G_g\bigl(p)-p\bigr)$. Throughout the section, fix $\lambda \ge \Delta^{-1/2}$, and define the product metric
\begin{equation}\label{eq:product-metric}
d_\lambda\bigl((p, g), (q, g')\bigr) := \ell^\Theta_{\mathrm C, \infty}(p, q) + \lambda \lvert g- g'\rvert.
\end{equation}
We will now be simultaneously using a coupled SKM recursions with the same step size, one for the categorical coefficients, and the other for the gain:
\begin{equation}\label{eq:coupled-skm}
\begin{gathered}
p_{k+1} = p_k + \alpha_k\bigl(H_{g_k}(p_k, Y_k) - p_k\bigr),\\
g_{k+1} = g_k + \alpha_k(R_k - g_k),
\end{gathered}
\end{equation}
where $Y_k = (S_k, R_k, S_{k+1})$. Intuitively, the coupled recursion can be seen as keeping an online estimate of the gain through an exponential moving average, and using that estimate to actively center sampled rewards. We now show that the one-sample augmented update is non-expansive in $d_\lambda$.
\begin{proposition}\label{prop:uncentered-nonexpansion}
For every uncentered sample $y = (s, r, s') \in \mathcal Y$, and all $(p, g), (q, g') \in \mathcal Z$,
\begin{equation}
d_\lambda\bigl( (H_g(p, y), r), (H_{g'}(q, y), r)\bigr) \le d_\lambda\bigl((p, g), (q, g')\bigr).
\end{equation}
\end{proposition}

The next proposition identifies the mean-field map associated with the stationary distribution $\mu$. 

\begin{proposition}\label{prop:mean-field-online}
Let $Y= (S, R, S')$ be a one-step sample from the stationary distribution of the policy-induced chain, i.e.,
\begin{equation}
S \sim \mu,\qquad \mathbb P(S'=j \mid S=i) = P_{ij},\qquad R \mid (S=i, S'=j) \sim \mathrm{Law}(R_{ij})
\end{equation}
for all $i, j \in \X$. Then
\begin{equation}
\E\bigl[\bigl(H_g(p, Y), R\bigr)\bigr] = \bigl(h^{(g)}_\mu(p), \bar r^\pi\bigr)\quad \text{for all } (p,g) \in \mathcal Z.
\end{equation}
Moreover, the map $(p, g) \mapsto\bigl(h^{(g)}_\mu(p), \bar r^\pi \bigr)$ is non-expansive in $d_\lambda$, and $\mathrm{fix}\bigl((p, g) \mapsto\bigl(h^{(g)}_\mu(p), \bar r^\pi \bigr)\bigr) = \{(p, \bar r^\pi): p \in \mathrm{fix}(\mathcal G)\}$.
\end{proposition}
\begin{corollary}\label{corol:residual-conv-online}
For every $p, g \in \mathcal Z$,
\begin{equation}
\mu_{\min}\ell^\Theta_{\mathrm C, \infty}\bigl(p, \mathcal G(p)\bigr) \le d_\lambda \bigl(\bigl(p, g\bigr), \bigl(h^{(g)}_\mu(p), \bar r^\pi\bigr) \bigr)\quad\text{and}\quad \lvert g- \bar r^\pi \rvert\le \lambda^{-1} d_\lambda \bigl(\bigl(p, g\bigr), \bigl(h^{(g)}_\mu(p), \bar r^\pi\bigr) \bigr).
\end{equation}
\end{corollary}

We now establish the non-asymptotic convergence of the coupled SKM recursion \eqref{eq:coupled-skm}. 
\begin{theorem}[Coupled recursion Markovian convergence]\label{thm:online-main}
Let $(S_k)_{k\ge 0}$ be an irreducible and aperiodic Markov chain with stationary distribution $\mu$, let $Y_k = (S_k, R_k, S_{k+1})$ be the associated uncentered one-sample process, and let $\lambda \ge \Delta^{-1/2}$. First, if $\alpha_k = (k+1)^{-a}$ with $\frac{4}{5} < a \le 1$, then the coupled SKM recursion \eqref{eq:coupled-skm} converges a.s.\ to a random limit $(p^\star, \bar r^\pi)$ such that $p^\star \in \mathrm{fix}(\mathcal G)$. Second, fix any $a_1 \in \bigl(\frac{4}{5}, 1\bigr)$ and set $\varepsilon_{a_1}:=(5a_1-4)/10$. There is an explicit two-phase step size schedule and constants $C^{\mathcal Z}_{\mathrm{mk}}, \kappa_{a_1} >0$ such that, for all $k \ge 1$, 
\begin{equation}
\begin{gathered}
\E\bigl[\ell^\Theta_{\mathrm C, \infty}\bigl(p_k, \mathcal G(p_k)\bigr)\bigr] \le \frac{C^{\mathcal Z}_{\mathrm{mk}}}{\mu_{\min}}(k+1)^{-1/10}\min \left\{\kappa_{a_1} \log (k+1), (k+1)^{\varepsilon_{a_1}} \right\},\\ 
\E[\lvert g_k - \bar r^\pi \rvert] \le \frac{C^{\mathcal Z}_{\mathrm {mk}}}{\lambda}(k+1)^{-1/10}\min \left\{\kappa_{a_1} \log (k+1), (k+1)^{\varepsilon_{a_1}} \right\},
\end{gathered}
\end{equation}
where $C_{\mathrm{mk}}^{\mathcal Z}$ depends on $a_1$, $\mathrm{diam}_{d_\lambda}(\mathcal Z) = \sqrt{\theta_d - \theta_1} + \lambda$, 
and on the corresponding finite-state Poisson constant for the uncentered chain $(Y_k)$. Similarly to Theorem~\ref{thm:centered-markov}, the fixed-exponent choice $a = \frac{4}{5}+\epsilon$ for  $\epsilon >0$ gives a.s.\ convergence with expected residual and gain error decaying at rate $O(k^{-1/10+\epsilon/2})$.

\begin{remark}
Appendix~\ref{sec:synchronous} shows that full synchronous sampled backups are gain-free at the exact quotient-law level: changing the centering constant only induces a common translation. However, this does not remove the need for online gain estimation in the practical fixed-grid categorical recursion, where projection can change the coefficients.
\end{remark}

\end{theorem}


\section{Experimental results}
We present two fixed-policy evaluation experiments: a simple five-state MDP to validate the theoretical claims of the paper, and a continuous-state experiment based on \texttt{Pendulum-v1} \citep{gym} in Appendix~\ref{sec:pendulum}, as a proof-of-concept of neural function approximation. 
The goal of the experiments is to empirically validate the theoretical results proposed in the paper regarding the quotient-categorical formulation and SA recursions. We use a fixed categorical support $\Theta$ with uniform spacing and evaluate all quantities using the supremum-Cramér metric in both of the experiments.

\noindent\textbf{Five-state MDP.}
We first consider a simple five-state MDP with deterministic rewards in $[0,1]$ and an irreducible and aperiodic transition matrix. Hence, the chain has a unique stationary distribution $\mu$, and the gain is given by $\bar r^\pi=\mu^\top r$. We compare four recursions: (i) the exact KM iteration for $\mathcal G$, (ii) the centered SKM recursion with i.i.d.\ one-step samples, (iii) the centered SKM recursion along a Markovian trajectory, and (iv) the coupled uncentered SKM recursion. For each method, we report the projected-operator residual $\ell^\Theta_{\mathrm C, \infty}(p_k,\mathcal G(p_k))$ and the corresponding mean-field residual $\ell^\Theta_{\mathrm C, \infty}(p_k,h_\rho(p_k))$. In the coupled case, we additionally report the gain error $|g_k-\bar r^\pi|$. The results in Figure~\ref{fig:mrp_main} match the predictions of the theory. In Figure~\ref{fig:mrp_main}(a), the exact KM residual $\ell^\Theta_{\mathrm C, \infty}(p_k,\mathcal G(p_k))$ decreases toward zero, illustrating the non-expansiveness and fixed-point structure of $\mathcal G$ established in Theorem~\ref{thm:gg-nonexp} and Corollary~\ref{corol:fixed-points}. The centered i.i.d.\ SKM recursion exhibits the same qualitative convergence behavior, conforming to Theorem~\ref{thm:centered-iid}. The plotted mean-field residual $\ell^\Theta_{\mathrm C, \infty}(p_k,h_\rho(p_k))$ is the residual controlled directly by the SA result, and Corollary~\ref{corol:residual-conversion} translates this control to the projected-operator residual $\ell^\Theta_{\mathrm C, \infty}(p_k,\mathcal G(p_k))$. The Markovian SKM curve decreases more slowly and shows a wider variability band, which is consistent with Theorem~\ref{thm:centered-markov}, where temporal dependence gives a weaker finite-time rate than in the i.i.d.\ setting. Figure~\ref{fig:mrp_main}(b) shows that the coupled gain-estimated recursion decreases both the categorical residual and the gain error $|g_k-\bar r|$, as predicted by Theorem~\ref{thm:online-main}. In particular, Corollary~\ref{corol:residual-conv-online} shows that the product residual controls both $\ell^\Theta_{\mathrm C, \infty}(p_k,\mathcal G(p_k))$ and $|g_k-\bar r|$. Finally, Figure~\ref{fig:mrp_main}(c) isolates the role of centering in the fixed-grid representation. When the fixed gain is set to $g=0$, the recursion targets $\mathcal G_0$ rather than the correctly-centered $\mathcal G$. Consequently, its residual with respect to $\mathcal G$ does not vanish, while the correctly-centered baseline continues to converge. This is consistent with the distinction between gain invariance at the quotient-law level and the dependence of fixed-grid categorical coefficients on the chosen representative.











\begin{figure}[t]
\centering
\includegraphics[width=\textwidth]{bellman_compatibility_final.png}
\caption{
\textbf{(a)} Centered exact and stochastic recursions validate the non-expansiveness and convergence predictions of Theorems~\ref{thm:gg-nonexp}, \ref{thm:centered-iid}, and \ref{thm:centered-markov}.
\textbf{(b)} Both the product residual and gain error decrease, validating Theorem~\ref{thm:online-main}.
\textbf{(c)} The wrong-centering ablation shows that using $g=0$ as the centering targets $\mathcal G_0$ rather than the correctly centered operator $\mathcal G$. The rate guide curves plot the theorem rates without their exact multiplicative constants, intended to show only the rate of decay.
}
\label{fig:mrp_main}
\end{figure}

\begin{figure}[t]
\centering
\includegraphics[width=\textwidth]{bellman_compatibility_distributions_prefi.png}
\caption{
Categorical bias laws learned in the five-state MRP.
\textbf{(a)} The exact KM fixed point gives the reference categorical fixed-point laws for the projected operator $\mathcal G$.
\textbf{(b)} The centered i.i.d.\ SKM recovers categorical laws closely aligned with the exact KM solution.
\textbf{(c)} The coupled gain-estimated SKM recovers a similar normalized solution while learning the gain online.
\textbf{(d)} The wrong-centering ablation produces a visibly degenerate fixed-grid solution, illustrating the effect of using $\mathcal G_0$ instead of $\mathcal G$.
Dashed vertical lines show the corresponding means.
}
\label{fig:mrp_distributions}
\end{figure}








Figure~\ref{fig:mrp_distributions} visualizes the learned quotient-categorical bias laws for each state in the five-state MRP. Each row corresponds to one state, and each panel shows the categorical law supported on the fixed grid $\Theta$, together with its mean indicated by a dashed vertical line. The exact KM fixed point, the centered i.i.d.\ SKM recursion, and the coupled SKM recursion recover closely aligned normalized categorical laws across states. Thus, after quotient normalization by a common shift, the stochastic recursions produce the same categorical fixed-point the exact projected iteration. This is consistent with Theorems~\ref{thm:centered-iid}, \ref{thm:centered-markov}, and~\ref{thm:online-main}, which establish convergence to $\mathrm{fix}(\mathcal G)$ under i.i.d.\ centered sampling, Markovian centered sampling, and Markovian online gain-estimated sampling, respectively.

In contrast, the wrong-centering ablation produces categorical laws that are visibly misaligned relative to the correctly centered solutions, even though the same fixed support $\Theta$ is used. This illustrates the distinction between quotient-level gain invariance and fixed-grid categorical effects. At the exact quotient-law level, changing the centering constant only induces a common translation, as Proposition~\ref{prop:tg} and Proposition~\ref{prop:sync-gain-free}. However, after projection onto a fixed categorical support, the coefficients depend on the chosen representative. Consequently, using an incorrect fixed gain $g=0$ takes the recursion toward the fixed point of $\mathcal G_0$ rather than the correctly centered operator $\mathcal G_{\bar r}$, which explains the shifted solution which has collapsed on the largest atom, presented in Figure~\ref{fig:mrp_distributions}.

\section{Conclusion}

We introduced a quotient-categorical formulation for average-reward distributional bias estimation, resolving the additive-constant ambiguity of the scalar bias at the level of state-indexed laws. This viewpoint leads to a projected operator that is well-defined and non-expansive in a coordinate Cram\'er metric, and to TD-style SA schemes with a.s.\ convergence and finite-iteration residual guarantees. When the gain is unknown, augmenting the categorical recursion with an online gain estimate preserves the same non-expansive structure and yields joint control of the projected residual and gain error. We empirically corroborated these predictions on a five-state MDP, and included a continuous-state neural function approximation experiment in Appendix~\ref{sec:pendulum}. The same machinery should also extend to multivariate categorical representations by replacing the one-dimensional cumulative embedding and projection with their multivariate analogues, following the construction of \citet{kaya2026finiteiterationtheoryasynchronouscategorical}. Our theory focuses on finite-state fixed-policy evaluation with fixed categorical supports. Theoretical analysis of richer function approximation and control settings is a natural next step.




\newpage
\bibliography{refs}

\newpage
\appendix
\section*{Appendix Table of Contents}
\addcontentsline{toc}{section}{Appendix}
\markboth{Appendix}{Appendix}
\startcontents[appendix]
\printcontents[appendix]{l}{1}{\setcounter{tocdepth}{3}}

\newpage
\section{Proofs of Section~\ref{sec:quotient}}\label{app:A}

\subsection{Proof of Proposition~\ref{prop:T}}
\begin{proof}[Proof of Proposition~\ref{prop:T}]
Fix $\eta \in \mathcal F^\X$, $c \in \R$, and $i \in \X$. For each $j$, convolution commutes with common translation:
\begin{equation}
\nu_{ij}\,\ast\,\bigl((\tau_c)_\# \eta_j \bigr) = (\tau_c)_\#(\nu_{ij}\,\ast\,\eta_j).
\end{equation}
Since pushforward by the same translation is linear on mixtures,
\begin{equation}
\begin{aligned}
\bigl(\mathcal T\bigl((\tau_c)_\#\eta\bigr)\bigr)_i &= \sum_{j \in \X}P_{ij} \,\nu_{ij}\,\ast\,\bigl((\tau_c)_\#\eta_j\bigr)\\
&=\sum_{j\in\X} P_{ij}\,(\tau_c)_\# (\nu_{ij}\,\ast\,\eta_j)\\
&=(\tau_c)_\# \sum_{j \in \X} P_{ij} (\nu_{ij}\,\ast\,\eta_j) = \bigl((\tau_c)_\# \mathcal T \eta \bigr)_i.
\end{aligned}
\end{equation}
This proves \eqref{eq:commute}. Therefore, if $\eta \approx \zeta$, then $\mathcal T \eta\approx \mathcal T \zeta$, so $\overline{\mathcal T}([\eta]) := [\mathcal T \eta]$ is independent of the representative.
\end{proof}

\subsection{Proof of Proposition~\ref{prop:tg}}
\begin{proof}[Proof of Proposition~\ref{prop:tg}]
Fix $i \in \X$. Since $R_{ij} - g' = (R_{ij} -g) + (g - g')$, we have
\begin{equation}
\nu^{(g')}_{ij} = (\tau_{g-g'})_\# \nu^{(g)}_{ij}.
\end{equation}
Using the shift-equivariance of translation and convolution:
\begin{equation}
\nu^{(g')}_{ij}\,\ast\,\eta_j = \bigl((\tau_{g-g'})_\#\nu^{(g)}_{ij} \bigr)\,\ast\,\eta_j = (\tau_{g-g'})_\#(\nu^{(g)}_{ij}\,\ast\,\eta_j).
\end{equation}
Summing over $j$ therefore yields
\begin{equation}
(\mathcal T_{g'}\eta)_i = (\tau_{g-g'})_\#(\mathcal T_g \eta)_i.
\end{equation}
Hence, $\mathcal T_{g'}\eta$ and $\mathcal T_g \eta$ differ by a common translation and thus determine the same equivalence class.
\end{proof}

\subsection{Proof of Lemma~\ref{lem:proj-shift-equiv}}
\begin{proof}[Proof of Lemma~\ref{lem:proj-shift-equiv}]
Write $\Pi^\Theta_{\mathrm C} \eta = \sum_{k=1}^d p_k \delta_{\theta_k}$. The coefficients $p_k$ are determined by the usual linear interpolation weights that split mass between neighboring support points in $\Theta$, together with edge clipping outside the interval $[\theta_1, \theta_d]$. Replacing $\eta$ by $(\tau_c)_\#\eta$ and the support $\Theta$ by $\Theta_c$ translates every relevant location by the same scalar $c$, so the interpolation weights are unchanged. Consequently, the projected coefficients are the same, and the projected law is translated by $c$.
\end{proof}

\subsection{Proof of Proposition~\ref{prop:cramer}}
\begin{proof}[Proof of Proposition~\ref{prop:cramer}]
Fix a state $i \in \X$ and a common shift $c \in \R$. Since both $\eta^{p, c}_i$ and $\eta^{q, c}_i$ are supported on the same ordered grid $\Theta_c$, the difference of their cdfs is constant on each interval $[\theta_k +c, \theta_{k+1} +c)$, and equals $\sum_{j=1}^k (p_{i, j} - q_{i,j})$. Therefore
\begin{equation}
\begin{aligned}
\ell_{\mathrm C}(\eta^{p, c}_i, \eta^{q, c}_i)^2 &= \sum_{k=1}^{d-1}(\theta_{k+1}-\theta_k)\Bigl(\sum_{j=1}^k(p_{i, j} - q_{i,j}) \Bigr)^2 \\
&=\Delta \sum_{k=1}^{d-1}\bigl(F_{p_i}(\theta_k) - F_{q_i}(\theta_k) \bigr)^2 = \ell^\Theta_{\mathrm C}(p_i, q_i)^2.
\end{aligned}
\end{equation}
Taking square roots proves the statewise identity, and the block-supremum identity follows by taking the maximum over the states.
\end{proof}

\subsection{Proof of Theorem~\ref{thm:gg-nonexp}}
We first state two lemmas.
\begin{lemma}\label{lem:applemma1}
For every $g \in \R$ and all $\eta, \zeta \in \mathcal F^\X$,
\begin{equation}
\ell_{\mathrm C, \infty}(\mathcal T_g\eta, \mathcal T_g\zeta) \le \ell_{\mathrm C, \infty}(\eta, \zeta).
\end{equation}
\end{lemma}
\begin{proof}
Fix $i \in \X$. The output law at state $i$ is a mixture of shifted successor laws, so
\begin{equation}
F_{(\mathcal T_g \eta)_i}(x) = \sum_{j\in\X} P_{ij}\E\bigl[F_{\eta_j}(x-
R_{ij}+g) \bigr],
\end{equation}
and similarly for $\zeta$. Jensen's inequality then gives
\begin{equation}
\ell_{\mathrm C}\bigl((\mathcal T_g\eta)_i, (\mathcal T_g\zeta)_i \bigr)^2 \le \sum_{j \in \X}P_{ij}\E\Bigl[\int_\R\bigl(F_{\eta_j}(x-R_{ij}+g) - F_{\zeta_j}(x-R_{ij}+g)\bigr)^2 \mathrm dx \Bigr].
\end{equation}
A change of variables gets rid of the translation, yielding
\begin{equation}
\ell_{\mathrm C}\bigl((\mathcal T_g\eta)_i, (\mathcal T_g\zeta)_i \bigr)^2 \le \sum_{j \in \X} P_{ij} \ell_{\mathrm C}(\eta_j, \zeta_j)^2 \le \ell_{\mathrm C, \infty}(\eta, \zeta)^2.
\end{equation}
Taking the maximum of the left-hand side over states and taking square roots proves the claim.
\end{proof}

\begin{lemma}\label{lem:applemma2}
For each $c \in \R$, the scalar categorical projection $\Pi^{\Theta_c}_\mathrm C$ is non-expansive in the Cram\'er metric. Consequently, its statewise product is non-expansive in $\ell_{\mathrm C, \infty}$.
\end{lemma}
\begin{proof}
This is the standard scalar categorical Cram\'er projection property, see \citet{rowland2018analysis}. Applying it statewise and taking the supremum over the blocks yield the claim.
\end{proof}
\begin{proof}[Proof of Theorem~\ref{thm:gg-nonexp}]
Fix $g \in \R$. To prove well-definedness, fix $p \in \Delta^\X_d$, and two representatives $\eta^{p, c}$ and $\eta^{p, c'}$. By Proposition~\ref{prop:tg},
\begin{equation}
\mathcal T_g \eta^{p, c'} = (\tau_{c'-c})_\# \mathcal T_g \eta^{p, c}.
\end{equation}
Applying Lemma~\ref{lem:proj-shift-equiv} statewise shows that the projected law family with shift $c'$ is the common-translation of the projected law family with shift $c$, so both representatives yield the same probability coefficients.

For non-expansiveness of $\mathcal G_g$ fix the representative shift $c = 0$. Then
\begin{equation}
\eta^{\mathcal G_g(p), 0} = \Pi^\Theta_\mathrm C(\mathcal T_g \eta^{p, 0}), \qquad \eta^{\mathcal G_g(q), 0} =\Pi^\Theta_\mathrm C(\mathcal T_g \eta^{q, 0}).
\end{equation}
Using Proposition~\ref{prop:cramer} and Lemmas \ref{lem:applemma1} and \ref{lem:applemma2},
\begin{equation}
\begin{aligned}
\ell^\Theta_{\mathrm C, \infty}\bigl(\mathcal G_g(p), \mathcal G_g(q)\bigr) &= \ell_{\mathrm C, \infty}(\eta^{\mathcal G_g(p), 0}, \eta^{\mathcal G_g(q), 0}) \\
&\le \ell_{\mathrm C, \infty}(\mathcal T_g \eta^{p,0}, \mathcal T_g \eta^{q, 0}) \\
&\le \ell_{\mathrm C, \infty}(\eta^{p, 0}, \eta^{q, 0}) = \ell^\Theta_{\mathrm C, \infty}\bigl(p, q).
\end{aligned}
\end{equation}
\end{proof}

\subsection{Proof of Corollary~\ref{corol:fixed-points}}
\begin{proof}[Proof of Corollary~\ref{corol:fixed-points}]
$\Delta^\X_d$ is a nonempty compact convex subset of a finite-dimensional Euclidean space, and $\mathcal G_g$ is continuous. Brouwer's fixed-point theorem \citep{aliprantis2006infinite} therefore guarantees the existence of fixed points.
\end{proof}

\newpage
\section{Proofs of Section~\ref{sec:centered}}\label{app:B}
For the convenience of proofs, let us write $L_b : \Delta_d \to \Delta_d$ for the unique coefficient map satisfying
\begin{equation}
\Pi^\Theta_\mathrm C \bigl((\tau_b)_\# \eta^{p,0} \bigr) = \eta^{L_b p,0},\qquad p \in \Delta_d, b\in\R.
\end{equation}

\subsection{Proof of Proposition~\ref{prop:g-block-form}}
\begin{proof}[Proof of Proposition~\ref{prop:g-block-form}]
Fix $g \in \R$ and $i \in \X$ and consider the representative shift $c=0$. By definition of $\mathcal G_g$,
\begin{equation}
\eta^{\mathcal G_g(p), 0}_i = \Pi^\Theta_\mathrm C\Bigl(\sum_{j \in \X} P_{ij}(\nu^{(g)}_{ij}\,\ast\,\eta^{p, 0}_j) \Bigr).
\end{equation}
Because scalar categorical projection is linear on mixtures,
\begin{equation}
\eta^{\mathcal G_g(p), 0}_i = \sum_{j \in \X}P_{ij}\Pi^\Theta_\mathrm C (\nu^{(g)}_{ij}\,\ast\,\eta^{p, 0}_j).
\end{equation}
Now, fix $j \in \X$ and let $B\sim \nu^{(g)}_{ij}$. Using the definition of $L_b$,
\begin{equation}
\Pi^\Theta_\mathrm C(\delta_b \,\ast\,\eta^{p, 0}_j) = \eta^{L_bp_j, 0}\qquad \text{for every } b \in \R.
\end{equation}
Taking expectation over $B$,
\begin{equation}
\Pi^\Theta_\mathrm C(\nu^{(g)}_{ij} \,\ast\,\eta^{p, 0}_j) = \E\bigl[\eta^{L_Bp_j,0} \bigr].
\end{equation}
Since each law inside the expectation is categorical on the same support $\Theta$, taking expectation simply averages the coefficient vectors, so
\begin{equation}
\E\bigl[\eta^{L_Bp_j,0} \bigr] = \eta^{\E[L_B p_j],0}.
\end{equation}
Comparing coefficients on the common support $\Theta$ shows that
\begin{equation}
\mathcal G_g(p)_i = \sum_{j \in \X} P_{ij} \E[L_B p_j], \qquad B \sim \nu^{(g)}_{ij}.
\end{equation}
\end{proof}

\subsection{Proof of Lemma~\ref{lem:one-sample-nonexp}}
\begin{proof}[Proof of Lemma~\ref{lem:one-sample-nonexp}]
Let $y=(s,b,s')$. For blocks $u \ne s$, the sampled update leaves the block unchanged. For the updated block,
\begin{equation}
\ell^\Theta_\mathrm C\bigl(H(p, y)_s, H(q, y)_s\bigr) = \ell^\Theta_\mathrm C(L_b p_{s'}, L_b q_{s'}).
\end{equation}
Using Proposition~\ref{prop:cramer} together with projection non-expansiveness and translation invariance of $\ell_\mathrm C$ gives
\begin{equation}
\ell^\Theta_\mathrm C(L_b p_{s'}, L_b q_{s'}) \le \ell^\Theta_\mathrm C(p_{s'}, q_{s'}).
\end{equation}
Taking the maximum over all state blocks proves the result.
\end{proof}

\subsection{Proof of Proposition~\ref{prop:mean-field-fix}}
\begin{proof}[Proof of Proposition~\ref{prop:mean-field-fix}]
The expectation identity follows blockwise from Proposition~\ref{prop:g-block-form}. $\mathrm{fix}(\mathcal G) \subseteq \mathrm{fix}(h_\rho)$ follows trivially from the definition of $h_\rho$. Conversely, if $h_\rho(p) = p$, then $\rho_i(\mathcal G(p)_i - p_i) = 0$ for every $i \in \X$. $\rho_{\min}>0$ implies that $\mathcal G(p)_i = p_i$ for every $i \in \X$, hence $p \in \mathrm{fix}(\mathcal G)$.

For the non-expansiveness statement, fix $i \in \X$. Since
\begin{equation}
h_\rho(p)_i = (1-\rho_i)p_i + \rho_i\mathcal G(p)_i,\qquad h_\rho(q)_i = (1-\rho_i)q_i + \rho_i \mathcal G(q)_i,
\end{equation}
we have for each $k = 1, \dots, d-1$,
\begin{equation}
F_{h_\rho(p)_i}(\theta_k) - F_{h_\rho(q)_i}(\theta_k) = (1 - \rho_i)\bigl(F_{p_i}(\theta_k) -F_{q_i}(\theta_k)\bigr) + \rho_i \bigl(F_{\mathcal G(p)_i}(\theta_k) -  F_{\mathcal G(q)_i}(\theta_k)\bigr).
\end{equation}
Through the triangle inequality with the vector of cumulative differences with $(d-1)$ elements,
\begin{equation}
\begin{aligned}
\ell^\Theta_\mathrm C\bigl(h_\rho(p)_i, h_\rho(q)_i \bigr) &\le (1-\rho_i) \ell^\Theta_\mathrm C(p_i, q_i) + \rho_i \ell^\Theta_\mathrm C\bigl(\mathcal G(p)_i, \mathcal G(q) _i\bigr) \\
&\le \ell^\Theta_{\mathrm C, \infty}(p, q),
\end{aligned}
\end{equation}
where the last inequality uses Theorem~\ref{thm:gg-nonexp}. Taking the maximum over $i \in \X$ proves non-expansiveness.
\end{proof}

\subsection{Proof of Corollary~\ref{corol:residual-conversion}}
\begin{proof}[Proof of Corollary~\ref{corol:residual-conversion}]
For each state $i \in \X$,
\begin{equation}
h_\rho(p)_i - p_i = \rho_i\bigl(\mathcal G(p)_i -p_i \bigr).
\end{equation}
Hence, for each $k = 1, \dots, d-1$,
\begin{equation}
F_{h_\rho(p)_i}(\theta_k) - F_{p_i}(\theta_k) = \rho_i\bigl(F_{\mathcal G(p)_i}(\theta_k) - F_{p_i}(\theta_k) \bigr),
\end{equation}
and therefore
\begin{equation}
\ell^\Theta_\mathrm C\bigl(p_i, h_\rho(p)_i\bigr) = \rho_i \ell^\Theta_\mathrm C\bigl(p_i, \mathcal G(p)_i\bigr).
\end{equation}
Taking the maximum over $i$ yields
\begin{equation}
\ell^\Theta_{\mathrm C, \infty}\bigl(p, h_\rho(p)\bigr) = \max_{i \in \X} \rho_i \ell^\Theta_\mathrm C\bigl(p_i, \mathcal G(p)_i\bigr),
\end{equation}
from which the stated bounds follow immediately.
\end{proof}

\subsection{Proof of Theorem~\ref{thm:centered-iid}}
We first state the following proposition.
\begin{proposition}[Two-phase i.i.d.\ residual refinement]\label{prop:centered-iid-two-phase}
Under the assumptions of Theorem~\ref{thm:centered-iid}, fix $a_1 \in \bigl(\frac23,1\bigr)$ and set
\begin{equation}
\varepsilon_{a_1}:=\frac{3a_1-2}{6},\qquad \kappa_{a_1}:=\frac{2^{\varepsilon_{a_1}}}{\log 2}.
\end{equation}
Let $T$ be the largest integer such that
\begin{equation}\label{eq:two-phase-threshold}
(k+1)^{\varepsilon_{a_1}}\le \kappa_{a_1}\log(k+1),\qquad 1\le k\le T.
\end{equation}
If the centered SKM recursion uses the two-phase step size
\begin{equation}\label{eq:two-phase-stepsize}
\alpha_k =
\begin{cases}
(k+1)^{-a_1}, & 1\le k\le T,\\
(k+1)^{-2/3}, & k>T,
\end{cases}
\end{equation}
then there exists $C_{\mathrm{iid}} >0$ such that, for all $k\ge1$,
\begin{equation}
\E\Bigl[\ell^\Theta_{\mathrm C,\infty}\bigl(p_k,h_\rho(p_k)\bigr)\Bigr]\le C_{\mathrm{iid}}(k+1)^{-1/6}\min\Bigl\{\kappa_{a_1}\log(k+1),(k+1)^{\varepsilon_{a_1}}\Bigr\}.
\end{equation}
Consequently,
\begin{equation}
\E\Bigl[\ell^\Theta_{\mathrm C,\infty}\bigl(p_k,\mathcal G(p_k)\bigr)\Bigr]\le \frac{C_{\mathrm{iid}}}{\rho_{\min}}(k+1)^{-1/6}\min\Bigl\{\kappa_{a_1}\log(k+1),(k+1)^{\varepsilon_{a_1}}\Bigr\}.
\end{equation}
\end{proposition}
\begin{proof}
Since the tail exponent in \eqref{eq:two-phase-stepsize} is the critical value $2/3$, the a.s.\ convergence assertion remains the separate $a>\frac23$ assertion in Theorem~\ref{thm:centered-iid}. The role of the two-phase schedule is to interpolate between the near-critical transient behavior of an exponent $a_1>\frac23$ and the critical last-iterate residual behavior.

We first record the constants used below. Let
\begin{equation}
D:=\sqrt{\theta_d-\theta_1},\qquad M_2:=\sqrt{|\X|}D,
\end{equation}
and define
\begin{equation}
S_{3a_1/2}:=\sum_{t=1}^{\infty}(t+1)^{-3a_1/2},\quad \omega_{a_1}:=\frac{1-2^{-a_1}}{2},\quad \nu_{a_1}:=\frac{1-2^{-a_1}}{1-a_1},\quad \eta_{a_1}:=1-(3/4)^{1-a_1}.
\end{equation}
Also let
\begin{equation}
R_{a_1}:=\frac{(4/3)^{a_1}}{1-4^{-a_1}},
\end{equation}
\begin{equation}
K_{a_1}:=\frac{D}{\sqrt{\pi\omega_{a_1}}} +\frac{2M_2}{\sqrt{\pi}}\left(\frac{S_{3a_1/2}}{\sqrt{\nu_{a_1}\eta_{a_1}}} +\frac{2^{1+a_1/2}R_{a_1}}{\sqrt{1-a_1}}\right),
\end{equation}
and
\begin{equation}
B_T:=\max\left\{\frac{2}{\sqrt3},\,8^{1/6}(T+2)^{\varepsilon_{a_1}}\right\}.
\end{equation}
The final constant can be chosen as
\begin{equation}\label{eq:two-phase-constant}
C_{\mathrm{iid}}:=\max\left\{K_{a_1}+6M_2,\; K_{a_1}B_T+\frac{2M_2\sqrt6}{\sqrt{\pi}}+6M_2\right\}.
\end{equation}

The threshold in \eqref{eq:two-phase-threshold} is well defined. Indeed, for $x\ge2$, the function $x^{\varepsilon_{a_1}}/\log x$ decreases until $x=\exp(1/\varepsilon_{a_1})$ and then increases to infinity. Its value at $x=2$ is $\kappa_{a_1}$. Hence the set of integers satisfying \eqref{eq:two-phase-threshold} is a nonempty initial interval, and after its last element $T$ the reverse inequality holds.

We now derive the residual estimate. Work in the cumulative-coordinate embedding associated with $\ell^\Theta_{\mathrm C,\infty}$. The block-sup Cram\'er metric is the block-sup norm of Euclidean cumulative-coordinate vectors, and this block-sup norm is bounded by the ambient Euclidean norm. Under the i.i.d.\ assumptions, the noise
\begin{equation}
U_k:=H(p_{k-1},Y_k)-h_\rho(p_{k-1})
\end{equation}
is a martingale difference and satisfies $\|U_k\|_{\mathrm C,2}\le M_2$. Define the averaged noise process
\begin{equation}
\overline U_k=(1-\alpha_k)\overline U_{k-1}+\alpha_k U_k,\qquad \overline U_0=0.
\end{equation}
By orthogonality of martingale differences in the Euclidean cumulative-coordinate space and Jensen's inequality,
\begin{equation}\label{eq:avg-noise-two-phase}
\E\bigl[\|\overline U_k\|_{\mathrm C,\infty}\bigr]\le \E\bigl[\|\overline U_k\|_{\mathrm C,2}\bigr] \le M_2\sqrt{\sum_{t=1}^k\alpha_t^2\prod_{r=t+1}^k(1-\alpha_r)^2}.
\end{equation}
Let
\begin{equation}
v_k:=\sum_{t=1}^k\alpha_t^2\prod_{r=t+1}^k(1-\alpha_r)^2.
\end{equation}
We claim that $v_k\le \alpha_k$ for all $k\ge1$. Since
\begin{equation}
v_k=(1-\alpha_k)^2v_{k-1}+\alpha_k^2,
\end{equation}
it is enough by induction to verify
\begin{equation}
\alpha_k^{-1}-\alpha_{k-1}^{-1}\le1.
\end{equation}
In the first phase this follows from the mean-value theorem:
\begin{equation}
(k+1)^{a_1}-k^{a_1}\le a_1 k^{a_1-1}<1.
\end{equation}
In the second phase,
\begin{equation}
(k+1)^{2/3}-k^{2/3}\le \frac23 k^{-1/3}<1.
\end{equation}
At the transition point,
\begin{equation}
\alpha_{T+1}^{-1}-\alpha_T^{-1} =(T+2)^{2/3}-(T+1)^{a_1} \le (T+1)^{2/3}-(T+1)^{a_1}+\frac23(T+1)^{-1/3}<1,
\end{equation}
because $a_1>2/3$. Therefore \eqref{eq:avg-noise-two-phase} gives
\begin{equation}\label{eq:avg-noise-final-two-phase}
\E\bigl[\|\overline U_k\|_{\mathrm C,\infty}\bigr]\le M_2\sqrt{\alpha_k}.
\end{equation}

Following the inexact Krasnosel'ski\u{\i}--Mann estimate of \citet[Theorem~2.10]{bravo2024stochastic}, with \eqref{eq:avg-noise-final-two-phase}, yields
\begin{equation}\label{eq:two-phase-core-bound}
\E\Bigl[\ell^\Theta_{\mathrm C,\infty}\bigl(p_k,h_\rho(p_k)\bigr)\Bigr]\le \frac{D}{\sqrt{\pi\tau_k}}+\frac{2M_2}{\sqrt{\pi}}\sum_{t=1}^{k-1}\frac{\alpha_t^{3/2}}{\sqrt{\tau_k-\tau_t}}+6M_2\sqrt{\alpha_k},
\end{equation}
where $\tau_k:=\sum_{t=1}^k\alpha_t(1-\alpha_t)$. It remains to bound the middle sum in the two phases.

\textbf{Primary phase.}
Assume $1\le k\le T$, so $\alpha_t=(t+1)^{-a_1}$. We use
\begin{equation}\label{eq:tau-primary-two-phase}
\omega_{a_1}(k+1)^{1-a_1}\le \tau_k\le \frac{1}{1-a_1}(k+1)^{1-a_1}.
\end{equation}
The upper bound is immediate by integral comparison. For the lower bound, $1-\alpha_t\ge1-2^{-a_1}$ and the standard estimate $\tau_k\ge \tau_1 k^{1-a_1}$ imply the displayed inequality.

Split the middle sum at $\lfloor k/2\rfloor$. For $t\le\lfloor k/2\rfloor$,
\begin{equation}
\tau_k-\tau_t\ge \nu_{a_1}\eta_{a_1}(k+1)^{1-a_1},
\end{equation}
and therefore
\begin{equation}
\sum_{t=1}^{\lfloor k/2\rfloor}\frac{\alpha_t^{3/2}}{\sqrt{\tau_k-\tau_t}} \le \frac{S_{3a_1/2}}{\sqrt{\nu_{a_1}\eta_{a_1}}}(k+1)^{-(1-a_1)/2}.
\end{equation}
For $t>\lfloor k/2\rfloor$,
\begin{equation}
\alpha_t^{1/2}\le 2^{a_1/2}(k+1)^{-a_1/2}.
\end{equation}
Using the left-shifted integral comparison
\begin{equation}
\frac{1}{\sqrt{\tau_k-\tau_t}}\le \frac{1}{\alpha_{t+1}(1-\alpha_{t+1})}\int_{\tau_t}^{\tau_{t+1}}\frac{\mathrm d\tau}{\sqrt{\tau_k-\tau}},
\end{equation}
and the ratio bound
\begin{equation}
\frac{\alpha_t}{\alpha_{t+1}(1-\alpha_{t+1})}\le R_{a_1},
\end{equation}
we get
\begin{equation}
\sum_{t=\lfloor k/2\rfloor+1}^{k-1}\frac{\alpha_t^{3/2}}{\sqrt{\tau_k-\tau_t}} \le \frac{2^{1+a_1/2}R_{a_1}}{\sqrt{1-a_1}}(k+1)^{-(1-a_1)/2}.
\end{equation}
Combining these estimates with \eqref{eq:two-phase-core-bound} and \eqref{eq:tau-primary-two-phase} gives
\begin{equation}\label{eq:primary-final-two-phase}
\E\Bigl[\ell^\Theta_{\mathrm C,\infty}\bigl(p_k,h_\rho(p_k)\bigr)\Bigr]\le (K_{a_1}+6M_2)(k+1)^{-(1-a_1)/2}.
\end{equation}
Since
\begin{equation}
(k+1)^{-(1-a_1)/2}=(k+1)^{-1/6}(k+1)^{\varepsilon_{a_1}},
\end{equation}
and \eqref{eq:two-phase-threshold} makes the minimum equal to $(k+1)^{\varepsilon_{a_1}}$ in this phase, \eqref{eq:primary-final-two-phase} proves the claimed bound for $k\le T$.

\textbf{Secondary phase.}
Now let $k>T$, so the tail stepsizes satisfy $\alpha_t=(t+1)^{-2/3}$. For any $t'\ge T$,
\begin{equation}\label{eq:tau-secondary-two-phase}
\tau_k-\tau_{t'}=\sum_{t=t'+1}^k\alpha_t(1-\alpha_t) \ge \frac12\sum_{t=t'+1}^k(t+1)^{-2/3}\ge \frac32\bigl((k+1)^{1/3}-(t'+1)^{1/3}\bigr).
\end{equation}
Split the middle sum in \eqref{eq:two-phase-core-bound} at $T$.

For the tail part $t>T$, we have $\alpha_t^{3/2}=(t+1)^{-1}$. By \eqref{eq:tau-secondary-two-phase},
\begin{equation}
\sum_{t=T+1}^{k-1}\frac{(t+1)^{-1}}{\sqrt{\tau_k-\tau_t}} \le \int_T^k\frac{(x+1)^{-1}}{\sqrt{\frac32((k+1)^{1/3}-(x+1)^{1/3})}}\,\mathrm dx.
\end{equation}
With $u=(x+1)^{1/3}$ and $c=(k+1)^{1/3}$, this integral equals
\begin{equation}
\sqrt6\int_{(T+1)^{1/3}}^{(k+1)^{1/3}}\frac{\mathrm d u}{u\sqrt{(k+1)^{1/3}-u}}
\end{equation}
and is bounded by
\begin{equation}\label{eq:tail-log-two-phase}
\sqrt6\,(k+1)^{-1/6}\log(k+1).
\end{equation}
To see the last step, write $x=((k+1)/(T+1))^{1/3}>1$ in the logarithmic expression obtained after integration:
\begin{equation}
\frac{\sqrt{x}+\sqrt{x-1}}{\sqrt{x}-\sqrt{x-1}} =2x-1+2\sqrt{x^2-x} \le 4x-1\le 3x^3\le (T+1)x^3=k+1.
\end{equation}

It remains to control the finitely many primary-phase terms $t\le T$. Evaluating the first two terms in \eqref{eq:two-phase-core-bound} at $T+1$ and using the primary-phase estimate gives
\begin{equation}
\frac{D}{\sqrt{\pi\tau_{T+1}}}+\frac{2M_2}{\sqrt{\pi}}\sum_{t=1}^{T}\frac{\alpha_t^{3/2}}{\sqrt{\tau_{T+1}-\tau_t}} \le K_{a_1}(T+2)^{-(1-a_1)/2}.
\end{equation}
If $k+1\ge 8(T+2)$, then \eqref{eq:tau-secondary-two-phase} with $t'=T+1$ gives
\begin{equation}
\tau_k-\tau_{T+1}\ge \frac34(k+1)^{1/3},
\end{equation}
and hence these early terms are bounded by
\begin{equation}
K_{a_1}\frac{2}{\sqrt3}(k+1)^{-1/6}.
\end{equation}
If instead $k+1<8(T+2)$, then
\begin{equation}
K_{a_1}(T+2)^{-(1-a_1)/2}\le K_{a_1}8^{1/6}(T+2)^{\varepsilon_{a_1}}(k+1)^{-1/6}.
\end{equation}
Thus, in both cases, the early terms are bounded by
\begin{equation}\label{eq:early-secondary-two-phase}
K_{a_1}B_T(k+1)^{-1/6}.
\end{equation}
Combining \eqref{eq:tail-log-two-phase}, \eqref{eq:early-secondary-two-phase}, and the terminal term $6M_2\sqrt{\alpha_k}=6M_2(k+1)^{-1/3}\le6M_2(k+1)^{-1/6}$, we obtain
\begin{equation}
\E\Bigl[\ell^\Theta_{\mathrm C,\infty}\bigl(p_k,h_\rho(p_k)\bigr)\Bigr] \le \left(K_{a_1}B_T+\frac{2M_2\sqrt6}{\sqrt{\pi}}+6M_2\right)(k+1)^{-1/6}\log(k+1).
\end{equation}
For $k>T$, the definition of the last crossing $T$ gives
\begin{equation}
(k+1)^{\varepsilon_{a_1}}>\kappa_{a_1}\log(k+1),
\end{equation}
so the minimum in the proposition is $\kappa_{a_1}\log(k+1)$. Since $\kappa_{a_1}\ge1$, the choice \eqref{eq:two-phase-constant} proves the residual bound in the secondary phase as well. The projected-operator residual follows from Corollary~\ref{corol:residual-conversion}.
\end{proof}
\begin{remark}
As a concrete example, $a_1=3/4$ gives $\varepsilon_{a_1}=1/24$, and the two-phase residual bound becomes
\begin{equation}
\E\Bigl[\ell^\Theta_{\mathrm C,\infty}\bigl(p_k,\mathcal G(p_k)\bigr)\Bigr]\le \frac{C_{\mathrm{iid}}}{\rho_{\min}}(k+1)^{-1/6} \min\left\{\frac{2^{1/24}}{\log2}\log(k+1),(k+1)^{1/24}\right\}.
\end{equation}
\end{remark}

We are now in a position to state the proof of Theorem~\ref{thm:centered-iid}.
\begin{proof}[Proof of Theorem~\ref{thm:centered-iid}]
Under the theorem assumptions, the mean-field map of the recursion is $h_\rho$ by Proposition~\ref{prop:mean-field-fix}. By Lemma~\ref{lem:one-sample-nonexp} the sampled map is samplewise $1$-Lipschitz, and by Proposition~\ref{prop:mean-field-fix} and Corollary~\ref{corol:fixed-points}, the mean-field map is non-expansive with nonempty fixed-point set. Compactness of $\Delta^\X_d$ implies bounded iterates and bounded martingale-difference noise. The almost sure convergence result of \citet[Theorem~2.5 and Example~2.7]{bravo2024stochastic} therefore applies to $h_\rho$ for every $\frac{2}{3}< a \le 1$.

For the expected residual, \citet[Theorem~4.4]{bravo2024stochastic} gives the following last-iterate profile for $\frac{1}{2}\le a \le 1$: there exist constants $C_1(a), C_2, C_3(a), C_4 >0$, depending only on $a$, the support diameter $\mathrm{diam}_{\ell^\Theta_{\mathrm C, \infty}}(\Delta^\X_d) = \sqrt{\theta_d - \theta_1}$, and any uniform bound on $\ell^\Theta_{\mathrm C, \infty}(H(p, y), p)$. Since both $H(p, y)$ and $p$ lie in $\Delta^\X_d$, this sampled-update term is itself bounded by the same diameter, so no additional problem-dependent constant is needed in our setting. Thus
\begin{equation}
\E\bigl[\ell^\Theta_{\mathrm C, \infty}\bigl(p_k, h_\rho(p_k)\bigr)\bigr] \le \begin{cases}
C_1(a)(k+1)^{-(a-\frac{1}{2})}, &\frac{1}{2}\le a < \frac{2}{3},\\
C_2\frac{\log(k+1)}{(k+1)^{1/6}}, &a=\frac{2}{3},\\
C_3(a)(k+1)^{-\frac{1-a}{2}}, &\frac{2}{3}<a<1,\\
C_4(\log(k+1))^{-\frac{1}{2}}, &a=1.
\end{cases}
\end{equation}
The two-phase residual assertion stated in the theorem is Proposition~\ref{prop:centered-iid-two-phase}. The projected-operator residual bound then follows from Corollary~\ref{corol:residual-conversion}.
\end{proof}

\subsection{Proof of Theorem~\ref{thm:centered-markov}}
We first state a lemma and then a proposition.
\begin{lemma}[Markovian inexact-KM reduction]\label{lem:markov-inexact-km}
Consider an SKM recursion on a compact subset of a finite-dimensional normed space, driven by a finite-state irreducible and aperiodic Markov chain, with samplewise non-expansive update map $H$ and non-expansive stationary mean-field map $h$. Let the step sizes be either $\alpha_k=(k+1)^{-a}$ with $a>\frac45$, or the matched two-phase schedule in Proposition~\ref{prop:centered-markov-two-phase}. Let
\begin{equation}
\tau_k:=\sum_{t=1}^k \alpha_t(1-\alpha_t),\qquad \sigma(0):=1,\qquad \sigma(u):=\min\{1,(\pi u)^{-1/2}\}\quad (u>0).
\end{equation}
There are constants $A_0,A_1,A_2>0$ such that the iterates satisfy
\begin{equation}\label{eq:markov-inexact-km}
\E\bigl[\|x_k-h(x_k)\|\bigr] \le A_0\sigma(\tau_k)\left(1+\sum_{t=2}^k \alpha_t\omega_{t-1}\right) +A_1\sum_{t=2}^k \alpha_t\sigma(\tau_k-\tau_t)\omega_{t-1} +A_1\omega_k,
\end{equation}
where the deterministic envelope $\omega_k$ obeys
\begin{equation}\label{eq:markov-poisson-envelope}
\omega_k\le A_2\tau_k\sqrt{\alpha_{k+1}}.
\end{equation}
\end{lemma}
\begin{proof}
This is the estimate obtained in the proof of \citet[Theorem~3.1]{blaser2026asymptotic} before specializing the last step to a single polynomial step size. The only Markovian ingredient is the Poisson-equation decomposition of
\begin{equation}
H(x_k,Y_{k+1})-h(x_k),
\end{equation}
which writes the Markovian noise as a martingale difference plus three telescoping error terms. Since the chain is finite and the update map is samplewise non-expansive on a compact set, the Poisson solution is uniformly bounded and Lipschitz on the compact state space. The auxiliary correction process $U_k$ in that proof therefore satisfies the deterministic envelope \eqref{eq:markov-poisson-envelope}. Applying the inexact Krasnosel'ski\u{\i}--Mann residual estimate to $z_k=x_k-U_k$ gives \eqref{eq:markov-inexact-km}; the three displayed terms correspond respectively to the initial-distance/inexactness constant, the convolution with the correction process, and the final correction $U_k$.

The polynomial form is used in \citet{blaser2026asymptotic} only when converting \eqref{eq:markov-inexact-km} into an explicit rate. The preceding Poisson-equation reduction uses the displayed regularity properties of the step sizes, namely monotonicity and bounded adjacent ratios, to control the telescoping error terms. These properties hold for the two-phase schedule because the transition between the two phases is continuous and each phase is polynomial. The critical-tail summations needed to turn \eqref{eq:markov-inexact-km} into a rate are verified explicitly in Proposition~\ref{prop:centered-markov-two-phase} below.
\end{proof}

\begin{proposition}[Two-phase Markovian residual refinement]\label{prop:centered-markov-two-phase}
Under the assumptions of Theorem~\ref{thm:centered-markov}, fix $a_1 \in \bigl(\frac45,1\bigr)$ and set
\begin{equation}
\varepsilon_{a_1}:=\frac{5a_1-4}{10},\qquad \kappa_{a_1}:=\frac{2^{\varepsilon_{a_1}}}{\log 2}.
\end{equation}
Let $T$ be the largest integer such that
\begin{equation}\label{eq:markov-two-phase-threshold}
(k+1)^{\varepsilon_{a_1}}\le \kappa_{a_1}\log(k+1),\qquad 1\le k\le T.
\end{equation}
Set $\gamma_T:=(T+1)^{-a_1}(T+2)^{4/5}$, and use the two-phase step size
\begin{equation}\label{eq:markov-two-phase-stepsize}
\alpha_k =
\begin{cases}
(k+1)^{-a_1}, & 1\le k\le T,\\
\gamma_T(k+1)^{-4/5}, & k>T.
\end{cases}
\end{equation}
Then there exists $C_{\mathrm{mk}} >0$ such that, for all $k\ge1$,
\begin{equation}
\E\Bigl[\ell^\Theta_{\mathrm C,\infty}\bigl(p_k,h_\mu(p_k)\bigr)\Bigr] \le C_{\mathrm{mk}}(k+1)^{-1/10} \min\Bigl\{\kappa_{a_1}\log(k+1),(k+1)^{\varepsilon_{a_1}}\Bigr\}.
\end{equation}
Consequently,
\begin{equation}
\E\Bigl[\ell^\Theta_{\mathrm C,\infty}\bigl(p_k,\mathcal G(p_k)\bigr)\Bigr] \le \frac{C_{\mathrm{mk}}}{\mu_{\min}}(k+1)^{-1/10} \min\Bigl\{\kappa_{a_1}\log(k+1),(k+1)^{\varepsilon_{a_1}}\Bigr\}.
\end{equation}
\end{proposition}
\begin{proof}
The threshold is well defined by the same argument as in Proposition~\ref{prop:centered-iid-two-phase}: the function $x^{\varepsilon_{a_1}}/\log x$ decreases from its value at $x=2$ until $x=\exp(1/\varepsilon_{a_1})$ and then increases to infinity. Thus the set of integers satisfying \eqref{eq:markov-two-phase-threshold} is a nonempty initial interval. The factor $\gamma_T$ makes the schedule continuous at the transition, since $\alpha_{T+1}=\alpha_T$, and hence non-increasing.

For $k\le T$, the recursion has only used the polynomial schedule $\alpha_t=(t+1)^{-a_1}$. The polynomial Markovian residual bound of \citet[Theorem~3.1]{blaser2026asymptotic} therefore gives
\begin{equation}
\E\Bigl[\ell^\Theta_{\mathrm C,\infty}\bigl(p_k,h_\mu(p_k)\bigr)\Bigr] \le C(k+1)^{-(1-a_1)/2}.
\end{equation}
Since
\begin{equation}
(k+1)^{-(1-a_1)/2} =(k+1)^{-1/10}(k+1)^{\varepsilon_{a_1}},
\end{equation}
and the minimum in the proposition is $(k+1)^{\varepsilon_{a_1}}$ in this phase, the desired bound follows for $k\le T$.

It remains to consider $k>T$. We apply Lemma~\ref{lem:markov-inexact-km} in the cumulative-coordinate norm that represents the block-sup Cram\'er metric. Since $T$ is fixed once $a_1$ is fixed, the two-phase schedule gives constants $0<c_\tau<C_\tau<\infty$ such that, for all $k>T$,
\begin{equation}\label{eq:markov-tau-tail}
c_\tau(k+1)^{1/5}\le \tau_k\le C_\tau(k+1)^{1/5}.
\end{equation}
Indeed, the tail contribution is a constant multiple of $\sum_{t>T}(t+1)^{-4/5}$, and the finitely many transition values are absorbed into $c_\tau$ and $C_\tau$. Combining \eqref{eq:markov-tau-tail} with \eqref{eq:markov-poisson-envelope} yields
\begin{equation}\label{eq:markov-omega-tail}
\omega_k\le C(k+1)^{-1/5}.
\end{equation}
Consequently,
\begin{equation}\label{eq:markov-first-error-sum}
\sum_{t=2}^k \alpha_t\omega_{t-1} \le C\left(1+\sum_{t=T+1}^k (t+1)^{-1}\right) \le C\log(k+1).
\end{equation}
The first term in \eqref{eq:markov-inexact-km} is therefore bounded by $C(k+1)^{-1/10}\log(k+1)$, and the final correction term \eqref{eq:markov-omega-tail} is smaller than this bound.

It remains to bound the convolution term in \eqref{eq:markov-inexact-km}. The finitely many terms with $t\le T$ contribute at most $C(k+1)^{-1/10}$ after increasing $C$, again using \eqref{eq:markov-tau-tail}; the finitely many near-transition cases are absorbed into the same constant. For the tail terms $T<t<k$, \eqref{eq:markov-omega-tail} gives $\alpha_t\omega_{t-1}\le C(t+1)^{-1}$, and the tail lower bound gives
\begin{equation}
\tau_k-\tau_t\ge c\bigl((k+1)^{1/5}-(t+1)^{1/5}\bigr).
\end{equation}
Thus
\begin{equation}
\sum_{t=T+1}^{k-1}\alpha_t\sigma(\tau_k-\tau_t)\omega_{t-1} \le C\sum_{t=T+1}^{k-1} \frac{(t+1)^{-1}}{\sqrt{(k+1)^{1/5}-(t+1)^{1/5}}}.
\end{equation}
By the same integral comparison used in the i.i.d.\ proof, with the change of variables $u=(x+1)^{1/5}$, the last display is at most
\begin{equation}
C\int_{(T+1)^{1/5}}^{(k+1)^{1/5}} \frac{\mathrm d u}{u\sqrt{(k+1)^{1/5}-u}} \le C(k+1)^{-1/10}\log(k+1).
\end{equation}
The omitted $t=k$ endpoint is bounded by $C(k+1)^{-1}$ and is therefore harmless. Combining these bounds in \eqref{eq:markov-inexact-km} gives
\begin{equation}
\E\Bigl[\ell^\Theta_{\mathrm C,\infty}\bigl(p_k,h_\mu(p_k)\bigr)\Bigr] \le C(k+1)^{-1/10}\log(k+1),\qquad k>T.
\end{equation}
For $k>T$, the definition of the largest crossing $T$ implies
\begin{equation}
(k+1)^{\varepsilon_{a_1}}> \kappa_{a_1}\log(k+1),
\end{equation}
so the minimum in the proposition is $\kappa_{a_1}\log(k+1)$. Since $\kappa_{a_1}\ge1$, enlarging the constant proves the claimed residual bound in the secondary phase. The projected-operator residual bound follows from Corollary~\ref{corol:residual-conversion}.
\end{proof}

We are now in a position to state the proof of Theorem~\ref{thm:centered-markov}.
\begin{proof}[Proof of Theorem~\ref{thm:centered-markov}] 
Because $\X$ and $\A$ are finite and the reward function $R:\X\times\A\to[0,1]$ is deterministic, the centered reward $B_k$ is finite even after conditioning only on $(S_k,S_{k+1})$. Hence the feasible centered sample set
\begin{equation}
\{(i, b,j) \in \X \times \R \times \X: \mathbb P(S_k =i, B_k = b, S_{k+1}=j) >0 \}
\end{equation}
is finite. The process $(Y_k)_{k\ge0}$ is Markov because the conditional law of $(B_{k+1}, S_{k+2})$ depends only on $S_{k+1}$, which is the last coordinate of $Y_k = (S_k, B_k, S_{k+1})$. It is irreducible and aperiodic on its feasible support because from any feasible triple $(i, b,j)$ to any other feasible triple $(i', b', j')$, one may first follow the irreducible and aperiodic state chain from $j$ to $i'$ and then realize the feasible transition producing $(b', j')$. By Lemma~\ref{lem:one-sample-nonexp}, the sampled map is samplewise non-expansive, and by Proposition~\ref{prop:mean-field-fix} and Corollary~\ref{corol:fixed-points}, the averaged operator is $h_\mu$, which is non-expansive with non-empty fixed-point set. The hypotheses of \citet{blaser2026asymptotic} are therefore satisfied, giving the fixed-exponent a.s.\ convergence assertion and the residual bound $C/\sqrt{\tau_k}$ with $\tau_k=\sum_{t=0}^{k-1}\alpha_t(1-\alpha_t)$. The two-phase residual assertion is Proposition~\ref{prop:centered-markov-two-phase}. The projected residual bounds once again follow from Corollary~\ref{corol:residual-conversion}.
\end{proof}

\subsection{Proof of Proposition~\ref{prop:counterexample}}
\begin{proof}[Proof of Proposition~\ref{prop:counterexample}]
Consider a one-state deterministic MDP. The transition matrix is $(1)$, the reward is almost surely $r$, and therefore $\bar r^\pi = r$. For any categorical state $p \in \Delta_d$, the centered update uses the shift $r - \bar r^\pi =0$, so $\mathcal G(p) = L_0 p = p$. Hence, $\mathcal G = \mathrm{Id}$. Now replace the centered reward by the uncentered reward in the one-state sampled update. The resulting update sends $p$ to $L_r p$. If the support has at least two atoms and $r \ne 0$, then $L_r p \ne p$ in general. For example if $p$ is a point mass at an interior atom, translating by $r$ and projecting back to the same support changes the cumulative coordinates. Therefore, the uncentered update is not an unbiased sample of $\mathcal G$, and in general it does not define an SA of the projected quotient operator.
\end{proof}

\newpage
\section{Proofs of Section~\ref{sec:uncentered}}\label{app:C}
For convenience of notation, let us write
\begin{equation}
\widehat H\bigl((p, g), (s, r, s') \bigr) := \bigl(H_g(p, (s, r, s')), r \bigr)
\end{equation}
and
\begin{equation}
\hat h_\mu(p, g) := \bigl(h^{(g)}_\mu(p), \bar r^\pi \bigr).
\end{equation}
This notation is more intuitive in making the gain estimate $g$ explicitly an argument of the recursion.
\begin{lemma}\label{lem:appC-1}
For every $b, b' \in \R$ and all $p, q \in \Delta_d$,
\begin{equation}
\ell^\Theta_\mathrm C(L_b p, L_{b'}q) \le \ell^\Theta_\mathrm C(p, q) + \Delta^{-1/2}\lvert b - b'\rvert.
\end{equation}
\end{lemma}
\begin{proof}
By the triangle inequality and the non-expansiveness of $L_{b'}$,
\begin{equation}
\ell^\Theta_{\mathrm C}(L_b p, L_{b'}q) \le \ell^\Theta_\mathrm C(L_b p, L_{b'}p) + \ell^\Theta_\mathrm C(L_{b'} p, L_{b'} q) \le \ell^\Theta_\mathrm C(L_b p, L_{b'}p) +\ell^\Theta_\mathrm C(p, q).
\end{equation}
It therefore suffices to bound $\ell^\Theta_\mathrm C(L_b p, L_{b'}p)$. Write $p = \sum_{m=1}^d p_m e_m$, where $e_m$ is the $m$th canonical basis vector. Since scalar categorical projection is linear on mixtures,
\begin{equation}
L_t p = \sum_{m=1}^d p_m L_t e_m\qquad\text{for every } t \in \R.
\end{equation}
Applying triangle inequality to the vector of cumulative masses yields
\begin{equation}\label{eq:sub-into-this}
\ell^\Theta_\mathrm C(L_b p, L_{b'}p) \le \sum_{m=1}^d p_m\ell^\Theta_\mathrm C(L_b e_m, L_{b'}e_m).
\end{equation}
Fix $m$ and define
\begin{equation}
v_m(t) := \bigl(F_{L_t e_m}(\theta_k)\bigr)_{k=1}^{d-1} \in \R^{d-1}.
\end{equation}
Since $\ell^\Theta_\mathrm C(u,v) = \sqrt{\Delta}\bigl\lVert (F_u(\theta_k) - F_v(\theta_k))_{k=1}^{d-1}\bigr\rVert_2$, it is enough to control $t \mapsto v_m(t)$. Let $x = \theta_{m+t}$ be the translated atom location. If $x \le \theta_1$, then $L_t e_m = e_1$ and $v_m(t) = \mathbf 1$. If $x \ ge \theta_d$, then $L_t e_m = e_d$ and $v_m(t) = 0$. If instead $x \in [\theta_k, \theta_{k+1}]$ for some $k \in \{1, \dots, d-1\}$, scalar categorical projection places mass
\begin{equation}
\lambda_k(x) := \frac{\theta_{k+1} -x}{\Delta}
\end{equation}
on $\theta_k$ and mass $1 - \lambda_k(x)$ on $\theta_{k+1}$. Hence
\begin{equation}
v_m(t) = \bigl(\underbrace{0,\dots,0}_{k-1 \text{ entries}},\lambda_k(x),\underbrace{1, \dots, 1}_{d-1-k\text{ entries}}\bigr),
\end{equation}
so on this interval $v_m$ is affine and
\begin{equation}
\frac{\mathrm d}{\mathrm dt} v_m(t) = -\frac{1}{\Delta}e_k.
\end{equation}
Therefore $\bigl\lVert\frac{\mathrm d}{\mathrm dt} v_m(t) \bigr\rVert_2 \le \Delta^{-1}$ wherever the derivative exists. Integrating this bound and multiplying by $\sqrt{\Delta}$ from the definition of $\ell^\Theta_\mathrm C$ gives
\begin{equation}
\ell^\Theta_\mathrm C(L_b e_m, L_{b'}e_m) \le \Delta^{-1/2}\lvert b- b'\rvert.
\end{equation}
Substituting this into \eqref{eq:sub-into-this} yields
\begin{equation}
\ell^\Theta_\mathrm C(L_b p, L_{b'}p) \le \Delta^{-1/2}\lvert b- b'\rvert
\end{equation}
\end{proof}

\begin{lemma}\label{lem:appC-2}
For all $g, g' \in \R$ and all $p, q \in \Delta^\X_d$,
\begin{equation}
\ell^\Theta_{\mathrm C, \infty}\bigl(\mathcal G_g(p), \mathcal G_{g'}(q) \bigr) \le \ell^\Theta_{\mathrm C, \infty}\bigl(p, q\bigr) + \Delta^{-1/2}\lvert g - g'\rvert.
\end{equation}
\end{lemma}

\begin{proof}
Fix $i \in \X$. By Proposition~\ref{prop:g-block-form},
\begin{equation}
\mathcal G_g(p)_i = \sum_{j\in \X} P_{ij} \E[L_{R_{ij} -g}p_j],\qquad \mathcal G_{g'}(q)_i = \sum_{j\in \X} P_{ij} \E[L_{R_{ij} -g'}q_j].
\end{equation}
For each $j$, let
\begin{equation}
u_j := \bigl( F_{L_{R_{ij} -g}p_j}(\theta_k) - F_{L_{R_{ij} -g'}q_j}(\theta_k)\bigr)^{d-1}_{k=1} \in \R^{d-1}.
\end{equation}
Since cumulative masses are linear in the coefficient vector, using triangle inequality,
\begin{equation}
\begin{aligned}
\ell^\Theta_\mathrm C\bigl( \mathcal G_g(p)_i, \mathcal G_{g'}(q)_i\bigr) &\le \sqrt{\Delta} \bigl\lVert\sum_{j \in \X}P_{ij}\E[u_j]\bigr\rVert_2 \\
&\le \sum_{j \in \X} P_{ij} \E \bigl[\ell^\Theta_\mathrm C(L_{R_{ij} - g}p_j, L_{R_{ij}-g'}q_j) \bigr] \\
&\le \sum_{j\in \X} P_{ij} \E\bigl[\ell^\Theta_\mathrm C(p_j, q_j) + \Delta^{-1/2}\lvert g- g' \rvert \bigr] \\
&\le \ell^\Theta_{\mathrm C, \infty}(p, q) + \Delta^{-1/2}\lvert g - g'\rvert.
\end{aligned}
\end{equation}
Taking the maximum over $i \in \X$ proves the result.
\end{proof}

\subsection{Proof of Proposition~\ref{prop:uncentered-nonexpansion}}
\begin{proof}[Proof of Proposition~\ref{prop:uncentered-nonexpansion}]
Let $y = (s, r, s')$. The second coordinate of $\widehat H$ is the common reward sample $r$, so $\lambda\lvert r - r \rvert = 0$. For the first coordinate, only the block $s$ is updated, and Lemma~\ref{lem:appC-1} gives
\begin{equation}
\ell^\Theta_\mathrm C\bigl(H_g(p, y)_s, H_{g'}(q, y)_s \bigr) = \ell^\Theta_\mathrm C(L_{r -g}p_{s'}, L_{r-g'}q_{s'}) \le \ell^\Theta_\mathrm C(p_{s'}, q_{s'}) + \Delta^{-1/2}\lvert g -g' \rvert.
\end{equation}
Taking the maximum over states,
\begin{equation}
\ell^\Theta_{\mathrm C, \infty}\bigl(H_g(p, y), H_{g'}(q,y) \bigr) \le \ell^\Theta_{\mathrm C, \infty}(p, q) + \Delta^{-1/2}\lvert g- g'\rvert.
\end{equation}
Since $\lambda \ge \Delta^{-1/2}$ by definition,
\begin{equation}
d_\lambda\bigl(\widehat H\bigl((p, g), y\bigr), \widehat H\bigl((q, g'), y\bigr)\bigr) \le \ell^\Theta_{\mathrm C, \infty}(p, q) + \Delta^{-1/2}\lvert g - g'\rvert \le d_\lambda \bigl((p,g),(q,g') \bigr).
\end{equation}
\end{proof}

\subsection{Proof of Proposition~\ref{prop:mean-field-online}}
\begin{proof}[Proof of Proposition~\ref{prop:mean-field-online}]
The expectation identity follows from Proposition~\ref{prop:g-block-form} applied under the stationary one-step law:
\begin{equation}
\E\bigl[H_g(p, Y) \bigr] = h^{(g)}_\mu(p),\qquad \E[R] = \bar r^\pi.
\end{equation}
For non-expansiveness, use the same cumulative-distribution argument as in the proof of Proposition~\ref{prop:mean-field-fix}. Since
\begin{equation}
h^{(g)}_\mu(p)_i = (1-\mu_i)p_i + \mu_i\mathcal G_g(p)_i,\qquad h^{(g')}_\mu(q)_i = (1-\mu_i)q_i + \mu_i\mathcal G_{g'}(q)_i,
\end{equation}
we have for each $k=1,\dots,d-1$,
\begin{equation}
F_{h^{(g)}_\mu(p)_i}(\theta_k) - F_{h^{(g')}_\mu(q)_i}(\theta_k) = (1-\mu_i)\bigl(F_{p_i}(\theta_k) - F_{q_i}(\theta_k)\bigr) + \mu_i \bigl(F_{\mathcal G_g(p)_i}(\theta_k) - F_{\mathcal G_{g'}(q)_i}(\theta_k)\bigr).
\end{equation}
Triangle inequality gives, for each $i \in \X$,
\begin{equation}
\begin{aligned}
\ell^\Theta_\mathrm C\bigl(h^{(g)}_\mu(p)_i, h^{(g')}_\mu(q)_i\bigr) &\le (1-\mu_i) \ell^\Theta_\mathrm C(p_i, q_i) + \mu_i \ell^\Theta_\mathrm C\bigl(\mathcal G_g(p)_i, \mathcal G_{g'}(q)_i \bigr) \\
&\le (1-\mu_i)\ell^\Theta_{\mathrm C, \infty}(p, q) + \mu_i \bigl(\ell^\Theta_{\mathrm C, \infty}(p,q) + \Delta^{-1/2}\lvert g -g'\rvert \bigr) \\
&\le \ell^\Theta_{\mathrm C, \infty}(p, q) + \Delta^{-1/2}\lvert g- g' \rvert.
\end{aligned}
\end{equation}
Taking the maximum over $i \in \X$ yields 
\begin{equation}
\ell^\Theta_{\mathrm C, \infty}\bigl(h^{(g)}_\mu(p), h^{(g')}_\mu(q) \bigr) \le \ell^\Theta_{\mathrm C, \infty}(p, q) + \Delta^{-1/2}\lvert g-g'\rvert,
\end{equation}
and therefore
\begin{equation}
d_\lambda\bigl(\hat h_\mu(p, g), \hat h_\mu(q, g')\bigr) \le \ell^\Theta_{\mathrm C, \infty}(p, q) + \Delta^{-1/2}\lvert g - g' \rvert \le d_\lambda\bigl((p, g), (q,g')\bigr).
\end{equation}
For the statement about fixed points, if $(p, g) \in \mathrm{fix}(\hat h_\mu)$, then necessarily $g = \bar r^\pi$. The first coordinate then reads $p = h_\mu(p)$, and $\mu_{\min} >0$ implies $p \in \mathrm{fix}(\mathcal G)$ by Proposition~\ref{prop:mean-field-fix}. The converse implication is immediate.
\end{proof}

\subsection{Proof of Corollary~\ref{corol:residual-conv-online}}
\begin{proof}[Proof of Corollary~\ref{corol:residual-conv-online}]
For $i \in \X$,
\begin{equation}
\mu_i \ell^\Theta_{\mathrm C}\bigl(p_i, \mathcal G_g(p)_i\bigr) = \ell^\Theta_{\mathrm C}\bigl(p_i, h^{(g)}_\mu (p)_i\bigr).
\end{equation}
Also,
\begin{equation}
\ell^\Theta_{\mathrm C}\bigl(\mathcal G_g(p)_i, \mathcal G(p)_i\bigr) \le \Delta^{-1/2}\lvert g - \bar r^\pi \rvert
\end{equation}
by Lemma~\ref{lem:appC-2} with $p = q$. Hence
\begin{equation}
\mu_i \ell^\Theta_{\mathrm C}\bigl(p_i, \mathcal G(p)_i\bigr) \le \ell^\Theta_{\mathrm C}\bigl(p_i, h^{(g)}_\mu(p)_i\bigr) + \Delta^{-1/2}\lvert g - \bar r ^\pi\rvert.
\end{equation}
Taking the maximum over $i \in \X$,
\begin{equation}
\mu_{\min}\ell^\Theta_{\mathrm C,\infty}\bigl(p, \mathcal G(p)\bigr) \le \ell^\Theta_{\mathrm C, \infty}\bigl(p, h^{(g)}_\mu(p)\bigr) + \Delta^{-1/2}\lvert g - \bar r^\pi\rvert \le d_\lambda\bigl((p, g), \hat h_\mu(p, g)\bigr),
\end{equation}
where the last step uses $\lambda \ge \Delta^{-1/2}$. The gain-error bound is immediate from the definition of $d_\lambda$.
\end{proof}

\subsection{Proof of Theorem~\ref{thm:online-main}}
\begin{proof}[Proof of Theorem~\ref{thm:online-main}]
Let $z_k = (p_k, g_k) \in \mathcal Z$. The coupled SKM recursion \eqref{eq:coupled-skm} can be written as
\begin{equation}
z_{k+1} = z_k + \alpha_k\bigl( \widehat H(z_k, Y_k) - z_k\bigr).
\end{equation}
View each categorical block $p_i$ through its cumulative masses $(F_{p_i}(\theta_1),\ldots,F_{p_i}(\theta_{d-1}))$, so that $\ell^\Theta_{\mathrm C,\infty}$ becomes the block-sup norm of these cumulative-coordinate vectors. On the augmented coordinates $(u,g)$, use the product norm $\lVert(u, g)\rVert_\lambda := \lVert u \rVert_{\mathrm{C},\infty} + \lambda \lvert g \rvert$. Under this identification, $d_\lambda$ is the induced distance on the compact convex set $\mathcal Z$. By Proposition~\ref{prop:uncentered-nonexpansion}, the sampled map is samplewise non-expansive in this norm. By Proposition~\ref{prop:mean-field-online} and Corollary~\ref{corol:fixed-points}, the mean-field map induced by the stationary one-step law is non-expansive with a non-empty set of fixed-points. Since rewards are finite-valued, the uncentered chain $(Y_k)_{k\ge0}$ is finite-state on its feasible support. The fixed-exponent Markovian theorem of \citet{blaser2026asymptotic} therefore applies. The final two bounds then follow from Corollary~\ref{corol:residual-conv-online}. The two-phase residual bound follows by applying Proposition~\ref{prop:centered-markov-two-phase}'s argument on the compact product space $(\mathcal Z, d_\lambda)$, with $\widehat H$ in place of $H$ and $\hat{h}_\mu$ in place of $h_\mu$. The only constants that change are the product-space diameter and the finite-state Poisson constant of the uncentered chain $(Y_k)_{k \ge 0}$. Corollary~\ref{corol:residual-conv-online} then gives the projected-operator and gain-error bounds.
\end{proof}

\newpage
\section{Synchronous gain-free structure}\label{sec:synchronous}

The previous section handles raw rewards by estimating the gain online. We close with a complementary exact-law observation: if one forms a full synchronous sampled backup before choosing categorical representatives, then the centering constant disappears already at the quotient level. This is a stronger, samplewise invariance statement than the fixed-point invariance of the exact quotient operator, but it lives in the idealized space of quotient distributions rather than in the practical fixed-grid categorical representation. To make the statement precise, let
\begin{equation}
\mathbf Y = \bigl((S'_i, R_i)\bigr)_{i \in \X}
\end{equation}
be a synchronous one-step sample such that, for each $i \in \X$,
\begin{equation}
\mathbb P(S'_i = j) = P_{ij},\qquad R_i \mid (S'_i=j) \sim \mathrm{Law}(R_{ij}),
\end{equation}
with the coordinates sampled independently across $i$. Thus $\mathbf Y$ contains one sampled successor and reward for every state. For $\eta \in \mathcal F^\X$ and $g \in \R$, define the exact synchronous sample map by
\begin{equation}
\bigl(\mathcal H_g(\eta, \mathbf Y)\bigr)_i := (\tau_{R_i -g})_\#\eta_{S'_i},\qquad i\in \X.
\end{equation}
Changing $g$ changes each centered reward $R_i-g$ by the same additive amount. Since the update is synchronous, that same shift appears in every output block.
\begin{proposition}\label{prop:sync-gain-free}
For every synchronous sample $\mathbf Y$, every $g, g' \in \R$, and every $\eta \in \mathcal F^\X$,
\begin{equation}
\mathcal H_{g'}(\eta, \mathbf Y) = (\tau_{g-g'})_\# \mathcal H_{g}(\eta, \mathbf Y).
\end{equation}
Consequently,
\begin{equation}
[\mathcal H_g(\eta, \mathbf Y)] = [\mathcal H_{g'}(\eta, \mathbf Y)]\quad \text{in } \mathcal F^\X / \approx.
\end{equation}
Thus the exact synchronous quotient sample is independent of the centering constant.
\end{proposition}
\begin{proof}
For every state $i$, $R_i - g' = (R_i - g) + (g-g')$, so $(\tau_{R_i - g'})_\# \eta_{S'_i} = (\tau_{g-g'})_\# \bigl((\tau_{R_i -g})_\# \eta_{S'_i} \bigr)$. The translation $(\tau_{g-g'})_\#$ is common to every state block, so the two output families represent the same element in the equivalence class.
\end{proof}
Taking the representative with $g=\bar r^\pi$ connects this gain-free quotient sample to the exact centered operator:
\begin{equation}
\E[\mathcal H_{\bar r^\pi}(\eta, \mathbf Y)] = \mathcal T \eta.
\end{equation}
To see this, condition on the successor state $S'_i = j$. The $i$th block of $\mathcal H_{\bar r ^\pi}(\eta, \mathbf Y)$  is then $(\tau_{R_i - \bar r^\pi})_\# \eta_j$ whose conditional law is $\nu_{ij}\,\ast\,\eta_j$. Averaging over $j$ with weights $P_{ij}$ gives exactly the $i$th component of $\mathcal T\eta$ from \eqref{eq:avg-operator}.

Thus synchronous exact samples provide a gain-free sample representation of the quotient-law target. This does not contradict the need for the online-gain categorical recursion above: once we choose a fixed categorical grid, projecting a translated representative can change the coefficients, so the exact quotient symmetry is not automatically preserved by a practical fixed-grid categorical update.

\newpage
\section{Neural function approximation results}\label{sec:pendulum}

\begin{figure}[!htb]
\centering
\begin{subfigure}{0.49\textwidth}
    \centering
    \includegraphics[width=\textwidth]{pendulum_exact_projected_residual_recap_final.png}
    \caption{Projected-operator residuals.}
    \label{fig:pendulum_projected_residual}
\end{subfigure}
\hfill
\begin{subfigure}{0.49\textwidth}
    \centering
    \includegraphics[width=\textwidth]{pendulum_product_residual_split_final.png}
    \caption{Product residual decomposition.}
    \label{fig:pendulum_product_residual}
\end{subfigure}
\caption{
Continuous-state \texttt{Pendulum-v1} evaluation.
\textbf{(a)} We evaluate the MC-trained critic, the online estimated gain critic, and the naive raw-reward critic against the correctly centered projected operator $\mathcal G$. Solid curves show the empirical supremum residual over validation states, and dotted curves show the residual averaged over states.
\textbf{(b)} For the online-gain critic, we report the product residual $d_\lambda$ together with its categorical and gain-estimation components. This corresponds to the decomposition used in Corollary~\ref{corol:residual-conv-online}.
}
\label{fig:pendulum_residuals}
\end{figure}

In this section, we consider a continuous-state fixed-policy evaluation problem based on \texttt{Pendulum-v1} \citep{gym}. This environment is a pendulum swing-up task with observation space $(\cos\theta,\sin\theta,\dot\theta)$ and a one dimensional bounded torque action $a\in[-2,2]$. The raw reward is
\begin{equation}
r_{\mathrm{raw}}(\theta,\dot\theta,a) = -\theta^2 - 0.1\dot\theta^2 - 0.001a^2,
\end{equation}
which penalizes angle error, angular velocity, and control effort. In the experiment we normalize this reward to $[0,1]$. The goal of this control problem is to keep the pendulum close to the upright position with small velocity and control cost. Since the state space is continuous, the tabular recursions from the previous experiment cannot be applied directly. Therefore, we use neural networks to parameterize the quotient-categorical critic. Data is collected using a fixed stochastic energy-based policy
\begin{equation}
a_t = \mathrm{clip}\!\left(K_E \dot\theta_t \cos\theta_t + K_D \dot\theta_t + \sigma\epsilon_t, -2,2 \right), \qquad \epsilon_t\sim\mathcal N(0,1).
\end{equation}
We compare the proposed gain-centered quotient-categorical critic with a naive categorical critic trained with raw rewards. We also train a scalar average-reward TD critic, which is used only as a reference for the learned categorical mean. Since the transition matrix is not available, the projected operator residual $\ell^\Theta_{C,\infty}(p_k,G(p_k))$ is estimated on held-out validation states using Monte Carlo (MC) one-step rollouts.

Figure~\ref{fig:pendulum_residuals} reports the main residual quantities in the continuous-state \texttt{Pendulum-v1} experiment. In Figure~\ref{fig:pendulum_residuals}(a), we evaluate the critic with MC-trained gain, the online-gain critic, and the raw-reward critic against the correctly-centered projected operator $\mathcal G$. The solid curves report the empirical supremum distance, corresponding to $\ell^\Theta_{\mathrm C,\infty}(p_k,\mathcal G(p_k))$, and the dotted curves report the averaged Cram\'er residual over states. Both proposed critics obtain substantially smaller residuals than the naive critic, which is consistent with Proposition~\ref{prop:counterexample}: the raw-reward categorical backup is not an SA of the correctly centered operator $\mathcal G$. Note that since this experiment is using neural function approximation, we don't expect the exact theoretical results to hold verbatim, and this experiment should be interpreted as a qualitative illustration of the arguments of the present work.

Figure~\ref{fig:pendulum_residuals}(b) evaluates the coupled online-gain recursion from Section~\ref{sec:uncentered}. Here, the critic is trained with the centered shift $r_t-g_t$, while the gain estimate $g_t$ is updated from the same reward samples. We report the product residual induced by the metric $d_\lambda$ in \eqref{eq:product-metric}, together with its categorical component and the gain component $\lambda |g_t-\bar r^\pi|$, where $\bar r^\pi$ is approximated by the MC gain estimate. The decrease of these terms is consistent with Theorem~\ref{thm:online-main}. Moreover, Corollary~\ref{corol:residual-conv-online} shows that this product residual controls both the projected-operator residual and the gain error. All of the experiments in this section were run on a cluster with four A100 GPUs over an estimated fifteen hours.
\end{document}
